%<*driver>
\def\nameofplainTeX{plain}
\ifx\fmtname\nameofplainTeX
\input l3docstrip.tex
\keepsilent
\askforoverwritefalse
\preamble

========================================================================
 osglecture-modes - generalized portable document modes

 (C) 2026 Matthias Werner
 E-mail: matthias.werner@informatik.tu-chemnitz.de

 Released under the LaTeX Project Public License v1.3c or later
 See https://www.latex-project.org/lppl.txt
========================================================================

\endpreamble
\postamble

This work is "maintained" (as per LPPL maintenance status) by
Matthias Werner.
\endpostamble

\generate{
  \file{osglecture-modes.sty}
    {\from{osglecture-modes.dtx}{package}}
  \file{osglecture-modes-overlays.code.tex}
    {\from{osglecture-modes.dtx}{overlays}}
  \file{osglecture-modes-ltxtalk.sty}
    {\from{osglecture-modes.dtx}{ltxtalk}}
}
% Lua cannot parse the TeX-style %% lines emitted by docstrip's default
% preamble and postamble.  Generate the helper separately without either.
\nopreamble
\nopostamble
\generate{
  \file{osglecture-modes-ltxtalk.lua}
    {\from{osglecture-modes.dtx}{lua}}
}
\expandafter\endbatchfile
\fi

\iffalse
%</driver>
%<*package>
\ProvidesExplPackage{osglecture-modes}
  {2026-08-02}
  {v0.4.0}
  {Portable document-modes}
%</package>
%<*driver>
\fi
\ExplSyntaxOff
\documentclass[add-bib=false]{osgdoc}

\usepackage{enumitem}
\usepackage{microtype}
\usepackage{needspace}
\newarg\modearg{<}{>}
\usepackage[
  languages={de,en},
  map={de=german,en=english},
  targetlang={job=3},
  load babel
]{langselect}
\usepackage{osglecture-modes}

\def\pkgname{osglecture-modes}
\def\thepkg{\pkg*{\pkgname}}

\makeindex
\setcnltx{
  name     = osglecture-modes,
  package  = osglecture-modes,
  version  = \UseVersionOf{osglecture-modes.sty},
  date     = \UseDateOf{osglecture-modes.sty},
  title    = {\ldeen{Portable Dokumentmodi}{Portable Document Modes}},
  info     = {\ldeen{Erweiterbare Modusgraphen für \LaTeX-Dokumente}
                    {Extensible mode graphs for \LaTeX\ documents}},
  authors  = {Matthias Werner[matthias.werner@informatik.tu-chemnitz.de]},
  abstract = {
    \ldeen{
      \LaTeX-Pakete wie \needpackage{beamer} und \needpackage{ltx-talk} bieten das 
      Konzept von \emph{Modi} an, bei dem, abhängig vom Dokumentmodus,
      nur bestimmte Dokumentteile gezielt übersetzt werden.

      Dieses Paket erweitert das Konzept und verwaltet Dokumentmodi 
      als gerichteten azyklischen Graphen. 
      Es ergänzt die Modusschnittstellen von \pkg*{beamer} und \pkg*{ltx-talk},
      funktioniert aber ebenso mit gewöhnlichen Dokumentklassen. Es gehört
      zum \pkg*{osglecture}-Bundle, ist jedoch auch unabhängig davon nutzbar.
    }{
       \LaTeX packages such as \needpackage{beamer} and \needpackage{ltx-talk} offer 
       the concept of \emph{modes}, in which, depending on the document mode, 
       only specific parts of the document are rendered.
      
      This package extends this concept and manages document modes as directed 
      acyclic graphs. It complements the mode interfaces of \pkg*{beamer}  and 
      \pkg*{ltx-talk}, but also works with standard document classes. 
      It is part of the \pkg*{osglecture} bundle, but can also be used independently of it.
    }
  },
  url = https://github.com/werner-matthias/osglecture,
  build-title
}
\OnlyDescription
\begin{document}

\section{\ldeen{Schnellstart}{Quick start}}

\ldeen{
Wer Beamer oder \pkg*{ltx-talk} verwendet, kennt Modusspezifikationen wie
\code{<presentation>} oder \code{<handout>}. \thepkg\ macht solche
Dokumentmodi auch in anderen Klassen nutzbar und erlaubt eigene Modi mit
mehreren Elternmodi. Ein gemeinsamer Quelltext kann so beispielsweise
Folien, ein Handout und ein Skript enthalten.
}{
Users of Beamer or \pkg*{ltx-talk} will recognize mode specifications such as
\code{<presentation>} or \code{<handout>}. \thepkg\ makes document modes
available in other classes as well and allows custom modes with multiple
parent modes. One source can thus contain slides, a handout, and a script.
}

\begin{example}[code-only,gobble=0]
\documentclass{article}
\usepackage[mode=script]{osglecture-modes}
\begin{document}
Common text.
\lecturemode<print>{Text for printed documents.}
\lecturemode<presentation>{Text for presentations.}
\IfLectureModeTF{script}{Detailed explanation.}{Short explanation.}
\end{document}
\end{example}

\ldeen{
Das Beispiel gibt den gemeinsamen Text, den Text für Druckfassungen und die
ausführliche Erklärung aus. Der Präsentationstext entfällt, weil
\code{script} den Modus \code{print}, aber nicht \code{presentation} aktiviert.
\code{mode=handout} würde stattdessen beide Textvarianten und die kurze
Erklärung auswählen. Die Paketoption ändert dabei nicht die Dokumentklasse.
}{
The example prints the common text, the text for printed documents, and the
detailed explanation. Presentation text is omitted: \code{script} activates
\code{print}, but not \code{presentation}. Selecting \code{mode=handout}
would print both variants and the short explanation instead. The package
option does not change the document class.
}

\ldeen{
Das Kernpaket benötigt einen aktuellen \LaTeX-Kernel. Es kann unabhängig
von \cls*{osglecture}, OLLM und Projektverzeichnissen verwendet werden.
Laden Sie es gewöhnlich nach der Dokumentklasse. Ohne explizite Moduswahl
erkennt es Beamer und \pkg*{ltx-talk}; bei anderen Klassen wählt es
\code{article}. Nur die portablen Textfilter aus
Abschnitt~\ref{sec:ltxtalk-filter} benötigen Lua\LaTeX.
}{
The core package needs a current \LaTeX\ kernel. It works independently of
\cls*{osglecture}, OLLM, and project directories. Normally, load it after the
document class. Without an explicit mode selection it detects Beamer and
\pkg*{ltx-talk}; for other classes it selects \code{article}. Only the
portable text filters in Section~\ref{sec:ltxtalk-filter} require Lua\LaTeX.
}

\section{\ldeen{Modi und Vererbung}{Modes and inheritance}}
\label{sec:model}

\ldeen{
Ein ausgewählter Modus aktiviert auch alle seine Eltern, deren Eltern usw.
Mehrere Eltern sind erlaubt; Kreise sind es nicht. Anders als bei einer
reinen Baumstruktur kann deshalb \code{handout} sowohl zu
\code{presentation} als auch zu \code{print} gehören. Die folgenden Modi
sind bereits deklariert; die Tabelle nennt jeweils die unmittelbaren Eltern.
}{
Selecting a mode also activates its parents, their parents, and so on.
Multiple parents are allowed, cycles are not. Unlike a simple tree,
this lets \code{handout} belong to both \code{presentation} and
\code{print}. The following modes are predefined; the table lists their
direct parents.
}

\begin{center}
\begin{tabular}{ll}
\textbf{\ldeen{Modus}{Mode}} & \textbf{\ldeen{Eltern}{Parents}} \\\hline
\code{presentation} & --- \\
\code{beamer}, \code{projector}, \code{trans} & \code{presentation} \\
\code{handout} & \code{presentation}, \code{print} \\
\code{longform}, \code{print} & --- \\
\code{article} & \code{longform}, \code{print} \\
\code{script} & \code{article} \\
\code{web} & \code{longform} \\
\code{*} & --- \\
\end{tabular}
\end{center}

\ldeen{
\code{slides} ist ein Alias für \code{projector}, kein zusätzlicher Modus.
\code{all} ist eine stets wahre Abfrage. Der Modus \code{*} dient der
Textfilterung (Abschnitt~\ref{sec:ltxtalk-filter}).
}{
\code{slides} is an alias for \code{projector}, not a separate mode.
\code{all} is an always-true query. Mode \code{*} is used for text filtering
(Section~\ref{sec:ltxtalk-filter}).
}

\subsection{\ldeen{Eigene Modi und mehrere aktive Modi}{Custom and simultaneous modes}}

\begin{example}[code-only,gobble=0]
\documentclass{article}
\usepackage{osglecture-modes}
\DeclareLectureMode[parents={article,print}]{reading}
\DeclareLectureModeAlias{reader}{reading}
\SetLectureModes{reader}
\FinalizeLectureModes
\begin{document}
\lecturemode<reading>{Reading edition.}
\lecturemode<longform>{Extended explanation.}
\end{document}
\end{example}

\ldeen{
\cs{SetLectureModes} ersetzt die bisherige Auswahl; \cs{AddLectureModes}
ergänzt sie. Beide nehmen kommaseparierte Listen entgegen. Die Auswahl
\code{slides,web} aktiviert beispielsweise gleichzeitig
\code{projector}, \code{presentation}, \code{web} und \code{longform}.
Die Dokumentklasse und ihre nativen Ausgabeoptionen bleiben davon unberührt.
}{
\cs{SetLectureModes} replaces the previous selection; \cs{AddLectureModes}
extends it. Both accept comma-separated lists. For instance,
selecting \code{slides,web} activates \code{projector},
\code{presentation}, \code{web}, and \code{longform} together. This does
not change the document class or its native output options.
}

\subsection{\ldeen{Zeitpunkt der Deklarationen}{When to declare modes}}
\ldeen{
Deklarieren Sie Modi, Eltern und Aliase und wählen Sie danach die aktiven
Modi, bevor die erste Abfrage erfolgt. \cs{FinalizeLectureModes} prüft den
gesamten Graphen, auch inaktive Modi, und schließt diese Konfiguration ab.
Die erste Abfrage mit \cs{IfLectureModeTF}, \cs{IfLectureModeT},
\cs{IfLectureModeF} oder einer spezifizierten \cs{lecturemode}-Anweisung
erledigt dies bei Bedarf automatisch; spätestens vor
\cs{begin}\Marg{document} wird abgeschlossen.
Danach sind Deklarationen, \cs{SetLectureModes} und \cs{AddLectureModes}
nicht mehr zulässig. Die Aktivierung bereits deklarierter Filtermodi im
Dokument ist ein eigener Vorgang, siehe Abschnitt~\ref{sec:ltxtalk-filter}.
}{
Declare modes, parents, and aliases and select the active modes before the
first query.

\cs{FinalizeLectureModes} validates the entire graph, including
inactive modes, and closes this configuration. The first query using
\cs{IfLectureModeTF}, \cs{IfLectureModeT}, \cs{IfLectureModeF}, or a
\cs{lecturemode} instruction with a specification does this automatically
if necessary; finalization takes place before \cs{begin}\Marg{document}
at the latest. Afterwards, further declarations and changes to the mode selection are
no longer allowed. Activating previously declared
filter modes in the document is a separate operation; see
Section~\ref{sec:ltxtalk-filter}.
}

\section{\ldeen{Inhalte und Werte auswählen}{Selecting content and values}}
\label{sec:selection}
\begin{commands}
\command{lecturemode}[\modearg{expression}\marg{content}]
\ldeen{Gibt den Inhalt aus, wenn der Ausdruck wahr ist. Ohne Spezifikation
wird der Inhalt immer ausgegeben.}{Prints the content if the expression is
true. Without a specification the content is always printed.}
\command{IfLectureModeTF}[\marg{expression}\marg{true}\marg{false}]
\ldeen{Führt je nach Ergebnis den zweiten oder dritten Parameter aus.}
{Executes the second or third argument according to the query result.}
\command{IfLectureModeT}[\marg{expression}\marg{true}]
\ldeen{Führt den zweiten Parameter nur bei wahrer Abfrage aus.}
{Executes the second argument only if the query is true.}
\command{IfLectureModeF}[\marg{expression}\marg{false}]
\ldeen{Führt den zweiten Parameter nur bei falscher Abfrage aus.}
{Executes the second argument only if the query is false.}
\end{commands}

\ldeen{
Diese Abfragen akzeptieren einen Modusnamen oder mehrere Namen, getrennt
durch \code{||} oder Kommas (Oder-Verknüpfung). Aliase sind erlaubt;
unbekannte Namen ergeben einen Fehler. \code{all} trifft immer zu.
Ein einzelnes \code{|}, Overlaynummern, Negation und Und-Verknüpfung gehören
nicht zu dieser Abfragesyntax. Für verschachtelte Bedingungen können Sie
mehrere Abfragen ineinandersetzen. Bei nativen Overlaybefehlen gelten andere
Regeln für Kommas und \code{|}, siehe Abschnitt~\ref{sec:backend-interaction}.
}{
These queries accept a mode name or several names separated by \code{||}
or commas (disjunction). Aliases are allowed; unknown names cause an error.
\code{all} always matches. A single \code{|}, overlay numbers, negation,
and conjunction are not part of this query syntax. Nest queries to express
combined conditions. Native overlay commands use different rules for commas
and \code{|}; see Section~\ref{sec:backend-interaction}.
}
\begin{example}[code-only,gobble=0]
\lecturemode<article||handout>{Material for readers.}
\IfLectureModeT{print}{\IfLectureModeF{handout}{Extended printed text.}}
\end{example}

\subsection{\ldeen{Modusabhängige Werte}{Mode-dependent values}}
\begin{commands}
\command{ModeValue}[\marg{specification}]
\ldeen{Expandiert zu einem ausgewählten Wert. Die Spezifikation besteht aus
mit \code{|} getrennten Einträgen der Form \code{modus=wert}.
\code{default=wert} oder ein Eintrag ohne \code{=} gibt den Ersatzwert an.
Ohne Treffer und ohne Ersatzwert ist das Ergebnis leer.}
{Expands to a selected value. The specification consists of entries of the
form \code{mode=value}, separated by \code{|}. Use \code{default=value}
or an entry without \code{=} for the fallback. With no match and no
fallback the result is empty.}
\end{commands}
\begin{example}[code-only,gobble=0]
% With graphicx loaded:
\includegraphics[width=\ModeValue{5cm|handout=8cm}]{plot}
\end{example}
\ldeen{
Die Reihenfolge der Einträge in der Spezifikation entscheidet nicht über
die Auswahl. Maßgeblich ist \cs{ActiveLectureModes}: Eltern stehen vor ihren
Kindern, die ausgewählten Modi werden in Auswahlreihenfolge durchlaufen;
jeder Modus kommt nur einmal vor. Der erste Modus mit einem Eintrag gewinnt,
\code{default} kommt zuletzt. Verwenden Sie hier die kanonischen Namen,
also beispielsweise \code{projector} statt des Alias \code{slides}.
}{
The order of entries in the specification does not determine the result.
\cs{ActiveLectureModes} supplies the priority: parents precede their
children, selected modes are traversed in selection order, and each mode
occurs only once. The first mode with an entry wins; \code{default} is
last. Use canonical names here, for example \code{projector} rather than
the alias \code{slides}.
}
\begin{example}[code-only,gobble=0]
% mode=script: ActiveLectureModes = longform,print,article,script
\ModeValue{script=8cm|print=6cm|default=5cm} % -> 6cm
\ModeValue{script=8cm|default=5cm}           % -> 8cm
\end{example}
\ldeen{
\cs{ModeValue} ist expandierbar und auch für Längen- und Registerzuweisungen
geeignet. Es finalisiert den Graphen nicht selbst: Soll es bereits in der
Präambel verwendet werden, rufen Sie zuvor \cs{FinalizeLectureModes} auf.
Werte dürfen weitere Gleichheitszeichen und Doppelpunkte enthalten;
\code{|} dient als Trenner. Verwenden Sie höchstens einen Ersatzwert.
}{
\cs{ModeValue} is expandable and can be used in length and register
assignments. It does not finalize the graph itself: call
\cs{FinalizeLectureModes} first when using it in the preamble. Values may
contain further equals signs and colons; \code{|} serves as the separator.
Use at most one fallback entry.
}

\clearpage
\section{\ldeen{Konfigurationsreferenz}{Configuration reference}}
\subsection{\ldeen{Paketoptionen}{Package options}}
\begin{options}
\keyval{mode}{name}\Default{auto}
\ldeen{Wählt beim Laden einen bereits bekannten Modus. Eigene Modi erst
nach dem Laden deklarieren und dann mit \cs{SetLectureModes} wählen.
Bei \code{auto} wählt Beamer im Handoutbetrieb \code{handout,beamer}, sonst
\code{slides,beamer}; \pkg*{ltx-talk} liefert entsprechend \code{handout},
\code{article} oder \code{slides}. Andere Klassen erhalten \code{article}.}
{Selects a mode already known at load time. Declare custom modes after
loading and select them with \cs{SetLectureModes}. With \code{auto}, Beamer
selects \code{handout,beamer} in handout mode and \code{slides,beamer}
otherwise; \pkg*{ltx-talk} selects \code{handout}, \code{article}, or
\code{slides} as appropriate. Other classes select \code{article}.}
\keybool{full}\Default{true}
\ldeen{Aktiviert die Integration in vorhandene Modus- und Overlaybefehle.
\code{full=false} lässt diese nativen Befehle unverändert. Portable Abfragen
bleiben nutzbar. Fehlende Befehle können weiterhin ergänzt werden.}
{Enables integration into existing mode and overlay commands.
\code{full=false} leaves these native commands unchanged. Portable queries
remain available. Missing commands may still be supplied.}
\keybool{patch-overlays}\Default{true}
\ldeen{Dieselbe Einstellung wie \option{full}; bei mehrfacher Angabe gilt
die zuletzt gesetzte Einstellung.}{The same setting as \option{full};
if specified more than once, the last setting wins.}
\keybool{ignorenonframetext}\Default{false}
\ldeen{Aktiviert zu Dokumentbeginn die Filterung mit \code{*}.
Bloße Angabe der Option bedeutet \code{true}.}
{Enables filtering with \code{*} at the start of the document.
Naming the option without a value means \code{true}.}
\keyval{raw-skip-envs}{list}
\ldeen{Ersetzt die Liste der beim Filtern wörtlich zu überspringenden
Umgebungen. Vorgabe: \env{verbatim}, \env{Verbatim}, \env{lstlisting}, \env{minted}, \env{comment}.
Siehe Abschnitt~\ref{sec:ltxtalk-filter}.}
{Replaces the list of environments to skip verbatim while filtering.
Default: \env{verbatim}, \env{Verbatim}, \env{lstlisting}, \env{minted}, \env{comment}.
See Section~\ref{sec:ltxtalk-filter}.}
\keybool{deferred}\Default{false}
\ldeen{Für Klassenautoren: unterdrückt die automatische Moduswahl und
Overlayintegration beim Laden. Die Klasse wählt den Modus selbst und ruft
nach dem Laden ihrer Basisklasse den unten beschriebenen
Aktivierungsbefehl auf.}
{For class authors: suppresses automatic mode selection and overlay
integration at load time. The class selects the mode itself and calls the
activation command described below after loading its base class.}
\end{options}

\subsection{\ldeen{Deklaration und Auswahl}{Declaration and selection}}
\begin{commands}
\command{DeclareLectureMode}[\oarg{parents=\meta{list}}\marg{mode}]
\ldeen{Deklariert einen Modus, optional mit kommaseparierten Eltern.
\code{parent} ist ein Synonym für \code{parents}. Eltern müssen spätestens
beim Abschluss deklariert sein; verwenden Sie ihre kanonischen Namen.}
{Declares a mode, optionally with comma-separated parents. \code{parent}
is a synonym for \code{parents}. Parents must be declared by finalization;
use their canonical names.}
\command{DeclareLectureModeParents}[\marg{mode}\marg{list}]
\ldeen{Setzt die vollständige Elternliste neu; ergänzt nicht die alte Liste.
Deklariert den Modus bei Bedarf mit.}
{Replaces the complete parent list rather than appending to it. Also
declares the mode if needed.}
\command{DeclareLectureModeAlias}[\marg{alias}\marg{mode}]
\ldeen{Deklariert einen alternativen Namen. Aliasketten sind erlaubt,
Aliaskreise und unbekannte Endziele nicht.}
{Declares an alternative name. Alias chains are allowed; alias cycles and
unknown final targets are not.}
\command{SetLectureModes}[\marg{list}]
\ldeen{Ersetzt die bisherige Auswahl durch die angegebene Liste.}
{Replaces the previous selection with the given list.}
\command{AddLectureModes}[\marg{list}]
\ldeen{Ergänzt die Auswahl. Die Modi müssen bereits
deklariert sein; Aliase werden aufgelöst, doppelte Einträge entfernt.}
{Extends the selection. Modes must already be declared;
aliases are resolved and duplicates removed.}
\command{FinalizeLectureModes}
\ldeen{Prüft und schließt den Graphen ab. Wiederholter Aufruf ist erlaubt.}
{Validates and finalizes the graph. Repeated calls are allowed.}
\command{LectureModes}
\ldeen{Expandiert zur kommaseparierten Auswahl ohne geerbte Eltern und mit
aufgelösten Aliasnamen.}{Expands to the comma-separated selection without
inherited parents and with aliases resolved.}
\command{ActiveLectureModes}
\ldeen{Expandiert zur aktiven Liste einschließlich Eltern. Vor Abschluss
ist diese Liste noch nicht berechnet. Beide Listenbefehle finalisieren
nicht selbst. Die Reihenfolge bestimmt auch \cs{ModeValue}.}
{Expands to the active list including parents. Before finalization this
list has not yet been computed. Neither list command finalizes the graph.
The order also determines \cs{ModeValue}.}
\command{LectureModesActivateOverlays}
\ldeen{Bindet nach dem Laden der Basisklasse die Overlaybefehle ein.
Nur bei \code{deferred=true} selbst einmal aufrufen. Wählt keinen Modus.}
{Integrates overlay commands after loading the base class. Call this once
only when using \code{deferred=true}. It does not select a mode.}
\end{commands}

\section{\ldeen{Zusammenspiel mit Backends}{Backend interaction}}
\label{sec:backend-interaction}

\ldeen{
  Besitzt die geladene Klasse keinen Befehl \cs{mode}, stellt das Paket dessen
  portable Inlineform bereit. Existiert bereits eine native Implementierung,
  behandelt das Paket ausschließlich Inlineausdrücke aus registrierten Modi.
  Ebenso sortieren die patchbaren Overlaybefehle und -Decoder von Beamer und
  \pkg*{ltx-talk} eine Winkelspezifikation in folgende Arten ein:
}{
  If the loaded class has no \cs{mode} command, the package provides its
  portable inline form. If a native implementation already exists, the package
  handles only inline expressions composed of registered modes. The patched
  overlay commands and decoders of Beamer and \pkg*{ltx-talk} likewise sort an
  angle-bracket specification into the following kinds:
}

\begin{itemize}
  \item \ldeen{
    \emph{portabel} -- eine mit \code{||} verknüpfte Oder-Liste registrierter
    Modi (oder \code{all}). Sie wird vollständig vom Paket ausgewertet und
    funktioniert auch ohne Overlaybackend.
  }{
    \emph{portable} -- a \code{||}-separated disjunction of registered modes
    (or \code{all}). Evaluated entirely by the package; works without any
    overlay backend.
  }
  \item \ldeen{
    \emph{modusqualifiziert} -- jedes mit \code{|} getrennte Segment ist ein
    registrierter Modus, eine Form \code{modus:rest} oder
    ein rein numerisches Overlayatom. Die portablen Anteile wertet das Paket
    aus; die aktiven nativen Reste werden zu einer einzigen
    Backendspezifikation zusammengesetzt. So wird \code{<slides:2-|handout:0>}
    portabel entschieden und dem Backend nur der zutreffende Overlaybereich
    gereicht.
  }{
    \emph{mode-qualified} -- every \code{|}-separated segment is a registered
    mode, a \code{mode:tail} form, or a purely numeric overlay
    atom. The package evaluates the portable parts and recombines the surviving
    native tails into a single backend specification. This lets
    \code{<slides:2-|handout:0>} be decided portably while only the matching
    overlay range reaches the backend.
  }
  \item \ldeen{
    \emph{nativ} -- alles Übrige: Backend-Modusnamen, \code{|}-Modusskopierung
    über unbekannte Modi, Actionspezifikationen und native Modusschalter mit unbekannten Namen. Wird unverändert an das Backend delegiert.
  }{
    \emph{native} -- everything else: backend mode names, \code{|} mode scoping
    over unknown modes, action specifications and native mode switches with unknown names. Delegated to the backend unchanged.
  }
\end{itemize}

\ldeen{
  Bei Backendbefehlen sind \code{|} und \code{,} innerhalb von \code{<...>} stets Backendsyntax;
  \code{||} ist der einzige portable Oder-Operator. Ein Ausdruck, der portable
  Modi mit nativen Overlayatomen mischt, muss dies also über \code{||} oder die
  Form \code{modus:rest} kenntlich machen. Fehlt der geladenen Klasse ein
  Overlaybackend, wird eine native oder unerkannte Spezifikation zu
  \emph{anzeigen} abgeschwächt. Alle diese Eingriffe lassen sich mit
  \code{patch-overlays=false} abschalten.
}{
  For backend commands, \code{|} and \code{,} inside \code{<...>} are always backend syntax; \code{||}
  is the only portable disjunction operator, so an expression that mixes
  portable modes with native overlay atoms must signal this with \code{||} or
  the \code{mode:tail} form. When the loaded class has no overlay backend, a
  native or unrecognised specification degrades to \emph{show}. All of these
  patches can be turned off with \code{patch-overlays=false}.
}

\subsection{\ldeen{Eigene Befehle overlayfähig machen}{Making custom commands overlay-aware}}
\begin{commands}
\command{MakeOverlayAwareCommand}[\oarg{options}\marg{command}]
\ldeen{Erweitert einen vorhandenen Befehl um eine unmittelbar folgende
Winkelspezifikation. Registrierte Modi werden portabel ausgewertet; echte
Overlays werden vom vorhandenen Backend ausgewertet.}
{Adds an angle-bracket specification immediately after an existing command.
Registered modes are evaluated portably; genuine overlays are evaluated
by the available backend.}
\end{commands}
\begin{options}
\keyval{arguments}{signature}\Default{\ldeen{leer}{empty}}
\ldeen{Argumentsignatur des Originalbefehls in der Syntax von
\cs{NewDocumentCommand}, beispielsweise \code{m} oder \code{s o}.
Sie muss alle nach der Winkelspezifikation folgenden Argumente beschreiben,
damit diese auch im ausgelassenen Zweig korrekt gelesen werden.}
{Argument signature of the original command in \cs{NewDocumentCommand}
syntax, for example \code{m} or \code{s o}. It must describe all arguments
following the angle specification so they are consumed correctly even when
the command is omitted.}
\keyval{default-mode}{mode}\Default{\ldeen{leer}{empty}}
\ldeen{Modusspezifikation für Aufrufe ohne Winkelargument. Ohne diese
Option wird ein Aufruf ohne Winkelargument unverändert ausgeführt.}
{Mode specification for calls without an angle argument. Without this
option, a call without an angle argument executes unchanged.}
\keychoice{fallback}{flatten,omit,error}\Default{flatten}
\ldeen{Verhalten bei echten Overlays ohne natives Backend:
\code{flatten} führt den Befehl aus, \code{omit} lässt ihn einschließlich
Argumenten aus, \code{error} meldet einen Fehler.}
{Behavior for genuine overlays without a native backend:
\code{flatten} executes the command, \code{omit} omits it and its arguments,
and \code{error} reports an error.}
\end{options}
\begin{example}[code-only,gobble=0]
\NewDocumentCommand{\Remark}{m}{\textbf{#1}}
\MakeOverlayAwareCommand[arguments=m]{\Remark}
% In the document:
\Remark<print>{For the printed edition.}
\end{example}
\ldeen{Die Winkelspezifikation steht vor Stern und optionalen Argumenten
 des Originalbefehls, beispielsweise:}
{The angle specification precedes the original command's star and optional
arguments, for example:}
\begin{example}[code-only,gobble=0]
\MakeOverlayAwareCommand
  [arguments={s o},default-mode=slides,fallback=flatten]{\\}
Text\\<slides>[3ex]
\end{example}

\section{\ldeen{Filtermodi}{Filter modes}}
\label{sec:ltxtalk-filter}

\ldeen{
  \thepkg\ verallgemeinert Beamers eigenes \cs{mode*} zu einem portablen
  Mechanismus: \code{*} ist selbst nur ein vorinstallierter Modus (siehe
  Abschnitt~\ref{sec:model}), und das Paket unterscheidet nicht zwischen
  ihm und jedem anderen, selbst deklarierten \emph{Filtermodus}. Ist ein
  Filtermodus aktiv, wird alles außerhalb registrierter
  \emph{Inhaltsumgebungen} (per Vorgabe \env{frame}, falls vorhanden)
  verworfen, mit registrierten Ausnahmen. Das funktioniert unter Beamer,
  \pkg*{ltx-talk} und mit gewöhnlichen Klassen wie \cls*{article}
  gleichermaßen -- eine Prüfung mit Aufgaben-, Lösungs- und
  Bewertungsumgebungen benötigt weder Beamer noch \pkg*{ltx-talk}.
}{
  \thepkg\ generalizes Beamer's own \cs{mode*} into a portable mechanism:
  \code{*} is itself just a pre-installed mode (see
  Section~\ref{sec:model}), and the package does not distinguish between
  it and any other, self-declared \emph{filter mode}. While a filter mode
  is active, everything outside registered \emph{unit environments} (by
  default \env{frame}, if present) is dropped, with registered
  exceptions. This works equally under Beamer, \pkg*{ltx-talk}, and with
  ordinary classes such as \cls*{article} -- an exam with task, solution,
  and grading environments needs neither Beamer nor \pkg*{ltx-talk}.
}

\ldeen{
  Die tatsächliche Filterung benötigt Lua\LaTeX\ und läuft, wo möglich,
  über das eigene Backend: Ist ein natives \cs{mode*} vorhanden (Beamer)
  \emph{und} sind keine Ausnahmen über dessen feste Positivliste hinaus
  registriert, delegiert \code{*} unverändert dorthin -- ohne Lua\LaTeX\
  und ohne Verhaltensänderung an bestehenden Beamer-Dokumenten. Sobald
  eigene Ausnahmen registriert sind oder ein anderer Filtermodus als
  \code{*} verwendet wird, übernimmt zwingend der portable Scanner dieses
  Pakets, da kein Backend eine beliebige, selbst deklarierte
  Positivliste kennen kann.
}{
  The actual filtering needs Lua\LaTeX\ and runs through the native
  backend where possible: if a native \cs{mode*} exists (Beamer)
  \emph{and} no exceptions beyond its fixed allow-list are registered,
  \code{*} delegates there unchanged -- no Lua\LaTeX, no behavior change
  for existing Beamer documents. As soon as custom exceptions are
  registered, or a filter mode other than \code{*} is used, the
  package's own portable scanner necessarily takes over, since no
  backend can know about an arbitrary, self-declared allow-list.
}

\begin{example}[code-only,gobble=0]
\usepackage[ignorenonframetext]{osglecture-modes}
\end{example}

\ldeen{
  Die Paketoption \code{ignorenonframetext} (Vorgabe bei bloßer Nennung:
  wahr) aktiviert den eingebauten Modus \code{*} -- entspricht also
  einem \cs{mode*} zu Dokumentbeginn. Für jeden
  anderen Filtermodus ist keine eigene Paketoption nötig: Er wird wie
  jeder andere Modus deklariert (\cs{DeclareLectureMode}) und über die
  Schalterform \cs{mode<name>} -- ohne Inhaltsargument -- direkt im
  Dokument aktiviert; \cs{mode*} ist dafür nur die aus Beamer-Kompatibilität
  beibehaltene Syntax für \cs{mode<*>}. \code{raw-skip-envs=\meta{Liste}}
  ersetzt die Liste wörtlich zu überspringender Umgebungen. Vorgabe sind
  \env{verbatim}, \env{Verbatim}, \env{lstlisting}, \env{minted} und
  \env{comment}. Ergänzen Sie eigene Verbatim-Umgebungen in dieser Liste,
  wenn sie in ausgefilterten Dokumentteilen vorkommen.

}{
  The package option \code{ignorenonframetext} (default when named bare:
  true) activates the built-in \code{*} mode -- equivalent to a
  \cs{mode*} at the start of the document. No package option is needed
  for any other filter mode: it is declared like any other mode
  (\cs{DeclareLectureMode}) and activated directly in the document via
  the switch form \cs{mode<name>} -- without a content argument;
  \cs{mode*} is only the syntax kept for Beamer compatibility with
  \cs{mode<*>}. \code{raw-skip-envs=\meta{list}} replaces the built-in
  list of environments to skip verbatim. The defaults are \env{verbatim},
  \env{Verbatim}, \env{lstlisting}, \env{minted}, and \env{comment}.
  Include custom verbatim environments in this list if they occur in
  filtered parts of the document.

}

\ldeen{
  Alles außerhalb einer registrierten Inhaltsumgebung wird verworfen, mit
  eingebauten Ausnahmen: \env{frame} und \env{frame*} selbst (falls
  vorhanden), die Gliederungsbefehle
  einschließlich \cs{appendix}, \cs{lecturemode}, \cs{mode} und \cs{maketitle}
  passieren den Scanner unverändert; die Gliederungsbefehle erhalten
  außerdem ein unmittelbar folgendes
  \cs{label}; ein \cs{label} an anderer Stelle verschwindet dagegen
  mit dem verworfenen Text. Der optionale, weiterhin
  \pkg*{ltx-talk}-spezifische Adapter \pkg*{osglecture-modes-ltxtalk}
  ergänzt \cs{lecture} als weitere allgemeine Ausnahme und setzt die
  Filterung nach jedem \cs{lecture} fort.
}{
  Everything outside a registered unit environment is dropped, with
  built-in exceptions: \env{frame} and \env{frame*} themselves (if
  present), structural commands
  including \cs{appendix}, \cs{lecturemode}, \cs{mode}, and
  \cs{maketitle} pass through the scanner unchanged; the structural
  commands also preserve an immediately following \cs{label}, while a \cs{label} elsewhere in
  dropped text vanishes along with it. The optional, still
  \pkg*{ltx-talk}-specific adapter \pkg*{osglecture-modes-ltxtalk} adds
  \cs{lecture} as a further exception and resumes filtering after every
  \cs{lecture}.
}

\begin{infobox}{\ldeen{Beginn der portablen Filterung}{When portable filtering starts}}
\ldeen{Im aktuellen Paket aktiviert der Modusschalter die Filterung,
das tatsächliche Überspringen beginnt aber erst nach \cs{maketitle}.
Beim Zusatzpaket für \pkg*{ltx-talk} startet es auch nach \cs{lecture}.
Text davor wird normal verarbeitet. Ein Dokument ohne einen solchen
Aufruf wird durch den Schalter allein noch nicht gefiltert. Dies betrifft
den portablen Filter; bei Delegation an Beamer gilt dessen natives Verhalten.}
{In the current package, the mode switch enables filtering, but actual
skipping starts only after \cs{maketitle}. With the additional package for
\pkg*{ltx-talk}, it also starts after \cs{lecture}. Earlier text is processed
normally. A document without such a call is not filtered by the switch
alone. This applies to the portable filter; delegation to Beamer retains
its native behavior.}
\end{infobox}

\subsection{\ldeen{Ausnahmen registrieren}{Registering exceptions}}
\begin{commands}
  \command{DeclareLectureKeepCommand}[\oarg{mode=\meta{name}}\marg{command}]
  \ldeen{Erhält einen bereits definierten Befehl mit genau einem
  Pflichtargument, etwa einen eigenen Eingabebefehl. Befehl und Argument
  werden normal ausgeführt.}{Keeps an already defined command with exactly
  one mandatory argument, such as a custom input command. The command and
  its argument are executed normally.}
  \command{DeclareLectureKeepLabelCommand}[\oarg{mode=\meta{name}}\marg{command}]
  \ldeen{Wie oben, aber nur unmittelbar nach einem Gliederungsbefehl.
  Damit kann ein eigener Label-Befehl das eingebaute \cs{label} ersetzen,
  ohne Labels in verworfenem Text zu erhalten.}
  {As above, but only immediately after a structural command. This lets a
  custom label command replace the built-in \cs{label} without preserving
  labels in discarded text.}
  \command{DeclareLectureKeepEnvironment}[\oarg{mode=\meta{name}}\marg{list}]
  \ldeen{Erhält die Umgebungen einer kommaseparierten Liste. Sie müssen
  bereits definiert sein. Ihr Inhalt wird normal verarbeitet; nach dem
  Ende einer solchen Umgebung wird wieder gefiltert.}
  {Keeps the environments in a comma-separated list. They must already be
  defined. Their content is processed normally; filtering resumes after
  each such environment ends.}
\end{commands}

\ldeen{
Ohne \code{mode=} gilt die Ausnahme in jedem Filtermodus. Mit
\code{mode=\meta{name}} gilt sie nur, wenn dieser Modus aktiv ist.
Dabei zählt auch die Vererbung: Ein ausgewählter Modus mit mehreren
Filtermodi als Eltern vereinigt deren Ausnahmen. Registrieren Sie die
Ausnahmen in der Präambel, bevor Sie die Filterung aktivieren.
}{
Without \code{mode=}, the exception applies in every filter mode. With
\code{mode=\meta{name}}, it applies only when that mode is active.
Inheritance counts too: a selected mode with multiple filter modes as
parents combines their exceptions. Register exceptions in the preamble
before activating filtering.
}

\begin{example}[code-only,gobble=0]
\documentclass{article}
\usepackage{osglecture-modes} % compile with LuaLaTeX
\NewDocumentEnvironment{task}{}{}{}
\DeclareLectureKeepEnvironment{task}
\DeclareLectureMode{solutions}
\NewDocumentEnvironment{solution}{}{}{}
\DeclareLectureKeepEnvironment[mode=solutions]{solution}
\title{Exercises}
\author{}
\date{}
\begin{document}
\mode<solutions>
\maketitle
This text is skipped.
\begin{task}An exercise.\end{task}
\begin{solution}Its solution.\end{solution}
\end{document}
\end{example}

\ldeen{
Das Beispiel gibt den Titel, die Aufgabe und die Lösung aus. Ersetzen Sie
\code{\string\mode<solutions>} durch \cs{mode*}, bleiben Titel und Aufgabe.
Ohne Aktivierung eines Filters wird auch der Text außerhalb der Umgebungen
normal verarbeitet; die Ausnahmeregistrierung allein blendet nichts aus.
Ein Schalter aktiviert den genannten Modus zusätzlich, er ersetzt nicht
die bisher aktive Modusmenge. Mehrere aktivierte Filtermodi vereinigen
daher ihre Ausnahmen. Eine portable Anweisung zum Deaktivieren bietet
das Paket derzeit nicht.
}{
The example prints the title, exercise, and solution. Replacing
\code{\string\mode<solutions>} with \cs{mode*} keeps the title and exercise.
Without an active filter, text outside the environments is processed
normally too; registering exceptions alone does not hide anything.
A switch activates the named mode additionally rather than replacing the
active mode set. Multiple activated filter modes therefore combine their
exceptions. The package currently provides no portable deactivation command.
}

\subsection{\ldeen{Zusatzpaket für ltx-talk}{Additional package for ltx-talk}}
\begin{example}[code-only,gobble=0]
\documentclass{ltx-talk}
\usepackage[ignorenonframetext]{osglecture-modes-ltxtalk}
\end{example}
\ldeen{
Das Zusatzpaket lädt das Kernpaket und ergänzt \cs{lecture} als Ausnahme.
Es akzeptiert \option{mode}, \option{ignorenonframetext} und
\option{raw-skip-envs}. Weitere Kernoptionen geben Sie bei Bedarf beim
vorherigen Laden von \thepkg\ an. Für eigene Umgebungen und Befehle gelten
dieselben Registrierungsbefehle wie oben.
}{
The additional package loads the core package and adds \cs{lecture} as an
exception. It accepts \option{mode}, \option{ignorenonframetext}, and
\option{raw-skip-envs}. If you need other core options, load \thepkg\ first
with those options. Register custom environments and commands as above.
}

\section{\ldeen{Modusabhängige Befehle definieren}{Defining mode-dependent commands}}
\label{sec:behavior}
\ldeen{
Für einfache Fallunterscheidungen genügt \cs{IfLectureModeTF}. Soll ein
öffentlicher Befehl je nach Dokumentmodus eine eigene Definition erhalten,
können Sie dagegen Verhaltensklassen verwenden. Diese bündeln Modi mit
derselben Befehlsdefinition, etwa \code{dynamic} für Folien und
\code{longform} für Artikel und Skripte. Verhaltensklassen und Dokumentmodi
sind getrennte Namenräume; gleiche Namen stellen keine automatische
Verbindung her. Das eigenständige Paket deklariert keine Verhaltensklassen.
}{
For simple conditional content, \cs{IfLectureModeTF} suffices. To give a
public command a different definition depending on the document mode,
use behavior classes instead. These group modes that share a command
definition, such as \code{dynamic} for slides and \code{longform} for
articles and scripts. Behavior classes and document modes have separate
namespaces; identical names do not establish an automatic connection.
The standalone package does not predeclare behavior classes.
}
\begin{example}[code-only,gobble=0]
\documentclass{article}
\usepackage[mode=script]{osglecture-modes}
\DeclareModeBehaviorClass{dynamic}
\DeclareModeBehaviorClass{longform}
\AssignLectureModeBehavior{slides}{dynamic}
\AssignLectureModeBehavior{script,article}{longform}
\NewDocumentCommand{\BriefRemark}{m}{\textbf{#1}}
\NewDocumentCommand{\FullRemark}{m}{\begin{quote}#1\end{quote}}
\DeclareModeAwareCommand{\Remark}{dynamic,longform}
\DeclareModeAwareCommandImplementation
  {\Remark}{dynamic}{\BriefRemark}
\DeclareModeAwareCommandImplementation
  {\Remark}{longform}{\FullRemark}
\begin{document}
\Remark{This is set as a quotation in script mode.}
\end{document}
\end{example}

\begin{commands}
\command{DeclareModeBehaviorClass}[\marg{class}]
\ldeen{Deklariert eine Verhaltensklasse.}{Declares a behavior class.}
\command{AssignLectureModeBehavior}[\marg{modes}\marg{class}]
\ldeen{Ordnet einer kommaseparierten Modusliste eine bereits deklarierte
Verhaltensklasse zu. Aliase sind erlaubt. Die Zuordnung wird nicht von
Elternmodi geerbt: Ordnen Sie die tatsächlich ausgewählten Modi zu.}
{Assigns a previously declared behavior class to a comma-separated list
of modes. Aliases are allowed. Assignment is not inherited from parent
modes: assign the modes actually selected.}
\command{DeclareModeAwareCommand}[\marg{command}\marg{classes}]
\ldeen{Deklariert einen neuen öffentlichen Befehl und die kommaseparierte
Liste der Verhaltensklassen, für die Definitionen vorliegen müssen.
Ein vorhandener Befehl wird nicht überschrieben.}
{Declares a new public command and the comma-separated list of behavior
classes for which definitions are required. An existing command is not
overwritten.}
\command{DeclareModeAwareCommandImplementation}
  [\marg{command}\marg{class}\marg{implementation command}]
\ldeen{Registriert den Definitionsbefehl für eine Verhaltensklasse.
Dieser besitzt selbst die vollständige Argumentsignatur. Verwenden Sie
für alle Varianten eine zum gemeinsamen Quelltext passende Signatur.}
{Registers the implementation command for a behavior class. That command
owns the complete argument signature. Use signatures compatible with the
shared source for all variants.}
\command{FinalizeModeAwareCommands}
\ldeen{Prüft alle geforderten Definitionen, auch für im aktuellen Lauf
inaktive Verhaltensklassen, und bindet die öffentlichen Befehle an ihre
gewählten Definitionen. Schließt bei Bedarf auch den Modusgraphen ab.}
{Checks all required definitions, including behavior classes inactive in
the current run, and binds public commands to their selected definitions.
Also finalizes the mode graph if necessary.}
\end{commands}

\ldeen{
Die Auswahl berücksichtigt die mit \cs{SetLectureModes} oder
\cs{AddLectureModes} gewählten Modi. Mindestens einer muss eine
Verhaltensklasse liefern; mehrere zugeordnete Modi dürfen keine
widersprüchlichen Klassen auswählen. Beim ersten Gebrauch eines solchen
Befehls, spätestens vor \cs{begin}\Marg{document}, erfolgt der Abschluss
automatisch. Danach sind keine weiteren Deklarationen oder Zuordnungen
für modusabhängige Befehle möglich. Später aktivierte Filtermodi wählen
die Befehlsdefinitionen nicht erneut aus.
}{
Selection considers the modes chosen with \cs{SetLectureModes} or
\cs{AddLectureModes}. At least one must supply a behavior class; multiple
assigned modes must not select conflicting classes. Finalization occurs
automatically on first use of such a command, or before
\cs{begin}\Marg{document} at the latest. Afterwards, no further declarations
or assignments for mode-dependent commands are possible. Activating filter
modes later does not reselect command definitions.
}
\section{\ldeen{Einschränkungen}{Limitations}}

\begin{itemize}
  \item \ldeen{
    Ein registrierter Modus darf einem nativen Overlayrest vorangestellt
    (\code{<slides:2->}), mit \code{|} modusweise gestaffelt
    (\code{<slides:2-|handout:0>}) oder per \code{||} oder-verknüpft werden. Ein
    Segment mit einem nicht registrierten Modusnamen macht jedoch die gesamte
    Spezifikation nativ, und eine Disjunktion mit nacktem \code{|} oder
    \code{,} über die Grenze zwischen portablem Modus und nativem Atom hinweg
    (\code{<presentation|2->}) ist nicht als portable Oder-Verknüpfung
    verfügbar; dafür ist \code{||} zu verwenden.
  }{
    A registered mode may prefix a native overlay tail (\code{<slides:2->}), be
    staggered per mode with \code{|} (\code{<slides:2-|handout:0>}), or be
    \code{||}-combined with one. A segment carrying an unregistered mode name
    makes the whole specification native, and a disjunction with a bare
    \code{|} or \code{,} across the boundary between a portable mode and a
    native atom (\code{<presentation|2->}) is not available as a portable
    \emph{or}; use \code{||} for that.
  }
  \item \ldeen{
    Laden Sie die Basisklasse vor \thepkg. Für Klassen, die erst später ihre
    Basisklasse laden, gibt es die Option \code{deferred=true}.
  }{
    Load the base class before \thepkg. Classes loading their base class
    later can use \code{deferred=true}.
  }
  \item \ldeen{
    Fehlt einer Basisklasse \cs{only}, \cs{uncover}, \cs{visible},
    \cs{onslide}, \cs{invisible}, \cs{alt}, \cs{pause} oder \cs{temporal},
    stellt \thepkg\ eine portable Fassung des \emph{Befehls} bereit. Die
    entsprechenden Environments (\env{onlyenv}, \env{uncoverenv},
    \env{visibleenv}, \env{invisibleenv}) bleiben dort undefiniert; gemeinsamer
    Quelltext sollte in diesem Fall die Befehlsform verwenden.
  }{
    If a base class lacks \cs{only}, \cs{uncover}, \cs{visible}, \cs{onslide},
    \cs{invisible}, \cs{alt}, \cs{pause}, or \cs{temporal}, \thepkg\ provides a
    portable version of the \emph{command}. The corresponding environments
    (\env{onlyenv}, \env{uncoverenv}, \env{visibleenv}, \env{invisibleenv})
    remain undefined there; shared source should use the command form in
    that case.
  }
\end{itemize}

\clearpage
\appendix
\section{\ldeen{Bundle-Integration: ein eigener Dokumenttyp}
{Bundle integration: a custom document type}}
\label{sec:bundle}
\ldeen{
Dieser Abschnitt richtet sich an Nutzer von \cls*{osglecture} und OLLM.
Für die eigenständige Verwendung von \thepkg\ sind diese Schritte nicht
nötig. Als Beispiel ergänzen wir \code{studyguide}, eine druckbare Langform
auf Basis von \cls*{scrbook}. Das OLLM-Target bestimmt den Dokumenttyp;
dieser wird als Modus ausgewählt. Das Dokumentprofil bestimmt dagegen
Basisklasse und Einrichtung des Dokuments.
}{
This section is for users of \cls*{osglecture} and OLLM. These steps are
not required for standalone use of \thepkg. We add \code{studyguide}, a
printable long-form document based on \cls*{scrbook}. The OLLM target
specifies the document type, which is selected as a mode. The document
profile specifies the base class and document setup instead.
}
\subsection{\ldeen{Target registrieren}{Registering the target}}
\ldeen{Legen Sie in einem eingebundenen OLLM-Definitionspfad die Datei
\code{targets/studyguide.toml} an:}
{Create \code{targets/studyguide.toml} in a configured OLLM definition path:}
\begin{example}[code-only,gobble=0]
schema = 1
kind = "target"
name = "studyguide"
version = "1.0"
doctype = "studyguide"
family = "publication"
unit_scopes = ["a", "as"]
\end{example}
\ldeen{Geben Sie das Target anschließend in \code{ollmconfig.toml} frei.
Der Definitionspfad muss über die lokale OLLM-Konfiguration oder ein
installiertes Bundle-Preset eingebunden sein.}
{Then enable the target in \code{ollmconfig.toml}. The definition path must
be configured through local OLLM settings or an installed bundle preset.}
\begin{example}[code-only,gobble=0]
[targets.studyguide]
languages = ["de", "en"]
\end{example}
\subsection{\ldeen{Modus und Verhalten festlegen}{Declaring mode and behavior}}
\ldeen{Erstellen Sie die für TeX auffindbare Datei
\code{osglecture-mode-studyguide.def}. Die Klasse \cls*{osglecture} stellt
die Verhaltensklasse \code{longform} bereits bereit.}
{Create \code{osglecture-mode-studyguide.def} where TeX can find it.
The \cls*{osglecture} class already supplies the \code{longform} behavior class.}
\begin{example}[code-only,gobble=0]
\ProvidesExplFile{osglecture-mode-studyguide.def}
  {2026-08-02}{v0.4.0}{studyguide mode setup}
\DeclareLectureMode[parents={longform,print}]{studyguide}
\AssignLectureModeBehavior{studyguide}{longform}
\end{example}
\ldeen{Damit sind Abfragen auf \code{studyguide}, \code{longform} und
\code{print} wahr. Soll ein Befehl anders arbeiten als in bisherigen
Langformen, deklarieren Sie eine eigene Verhaltensklasse und ergänzen Sie
die erforderlichen Befehlsdefinitionen nach Abschnitt~\ref{sec:behavior}.}
{Queries for \code{studyguide}, \code{longform}, and \code{print} now match.
If commands need behavior different from existing long-form documents,
declare a separate behavior class and supply the required command
definitions as described in Section~\ref{sec:behavior}.}
\Needspace{8\baselineskip}
\subsection{\ldeen{Dokumentprofil einrichten}{Setting up the document profile}}
\ldeen{Legen Sie \code{osglecture-profile-studyguide-scrbook.def} an:}
{Create \code{osglecture-profile-studyguide-scrbook.def}:}
\begin{example}[code-only,gobble=0]
\ProvidesExplFile{osglecture-profile-studyguide-scrbook.def}
  {2026-08-02}{v0.4.0}{studyguide scrbook profile}
\DeclareOsgLectureProfile{studyguide-scrbook}
  {
    backend = scrbook,
    class = scrbook,
    doctypes = {script,article,studyguide},
    mode-setup-file = osglecture-mode-studyguide.def,
    setup-file = osglecture-studyguide.code.tex
  }
\DeclareOsgLectureProfileClassOptions
  {studyguide-scrbook}{studyguide}{oneside}
\end{example}
\ldeen{
\code{mode-setup-file} enthält die Modusdeklarationen und Zuordnungen.
Die ebenfalls bereitzustellende Datei \code{osglecture-studyguide.code.tex}
enthält die Einrichtung und Formatierung nach dem Laden der Basisklasse.
Das Profil umfasst auch \code{script} und \code{article}, damit diese
Targets beim Ersetzen des projektweiten Langformprofils verfügbar bleiben.
Wählen Sie das Profil in der Projekt- oder Nutzerkonfiguration:
}{
\code{mode-setup-file} contains mode declarations and assignments. Also
provide \code{osglecture-studyguide.code.tex}, which contains setup and
formatting applied after loading the base class. Including \code{script}
and \code{article} keeps these targets available when replacing the
project's long-form profile. Select the profile in the project or user
configuration:
}
\begin{example}[code-only,gobble=0]
[latex.defaults]
script_profile = "studyguide-scrbook"
\end{example}
\ldeen{
Die Klasse lädt die Modusdatei vor Auswahl des neuen Modus. Eine Änderung
am Klassenkern ist dafür nicht nötig. Prüfen Sie abschließend den Build
für jede freigegebene Sprache: Target, Dokumenttyp und Profil müssen
zusammenpassen, die Modusabfragen das erwartete Ergebnis liefern und die
verwendeten modusabhängigen Befehle definiert sein. Prüfen Sie auch ein
Dokument mit explizitem \code{doctype} und \code{profile} ohne OLLM-Projekt.
Die vollständige Profil- und Targetkonfiguration beschreiben die
Dokumentationen von \cls*{osglecture} und OLLM.
}{
The class loads the mode setup file before selecting the new mode; no
change to the class core is needed. Finally, build each enabled language:
target, document type, and profile must agree, mode queries must return the
expected results, and the mode-dependent commands used must be defined.
Also check a document with explicit \code{doctype} and \code{profile}
outside an OLLM project. The \cls*{osglecture} and OLLM manuals describe
the full profile and target configuration.
}

% Comment out \OnlyDescription above to include the implementation.
\MaybeStop{}

\section{\ldeen{Implementierung}{Implementation}}
\subsection{\pkg*{osglecture-modes.sty}}
\AutoDocInput{osglecture-modes.dtx}{package}
\subsection{\pkg*{osglecture-modes-ltxtalk.sty}}
\AutoDocInput{osglecture-modes.dtx}{ltxtalk}
\subsection{\code{osglecture-modes-ltxtalk.lua}}
\AutoDocInput{osglecture-modes.dtx}{lua}

\end{document}

%</driver>
%<*package>
\prop_new:N \g_osglecture_modes_parents_prop
\prop_new:N \g_osglecture_modes_alias_prop
\prop_new:N \g_osglecture_modes_active_prop
\seq_new:N  \g_osglecture_modes_known_seq
\seq_new:N  \g_osglecture_modes_seed_seq
\seq_new:N  \g_osglecture_modes_active_order_seq
\bool_new:N \g_osglecture_modes_frozen_bool
\bool_new:N \g_osglecture_modes_patch_overlays_bool
\tl_new:N   \g_osglecture_modes_option_mode_tl

\msg_new:nnn { osglecture_modes } { frozen }
  { The~mode~graph~is~already~finalized;~'#1'~cannot~be~changed. }
\msg_new:nnn { osglecture_modes } { unknown-mode }
  { Unknown~mode~'#1'.~Declare~it~before~using~or~finalizing~the~graph. }
\msg_new:nnn { osglecture_modes } { unknown-parent }
  { Mode~'#1'~has~unknown~parent~'#2'. }
\msg_new:nnn { osglecture_modes } { unknown-alias-target }
  { Mode~alias~'#1'~has~unknown~target~'#2'. }
\msg_new:nnn { osglecture_modes } { alias-cycle }
  { Mode~alias~cycle~detected~at~'#1'. }
\msg_new:nnn { osglecture_modes } { graph-cycle }
  { Mode~parent~cycle~detected~at~'#1'. }

\bool_new:N \g_osglecture_modes_deferred_bool
\bool_new:N \g_osglecture_modes_ignorenonframetext_bool
\clist_new:N \g_osglecture_modes_raw_skip_envs_clist
\clist_gset:Nn \g_osglecture_modes_raw_skip_envs_clist
  {
    verbatim,
    Verbatim,
    lstlisting,
    minted,
    comment
  }

\keys_define:nn { osglecture_modes }
  {
    mode .tl_gset:N = \g_osglecture_modes_option_mode_tl,
    mode .initial:n = auto,
    patch-overlays .bool_gset:N =
      \g_osglecture_modes_patch_overlays_bool,
    patch-overlays .initial:n = true,
    full .meta:n = { patch-overlays = #1 },
    full .default:n = true,
    % \ldeen{Aktiviert den eingebauten Modus \code{*} (siehe die
    % Modell-Standardvorgaben weiter unten): Alles außerhalb einer
    % Unit-Umgebung (etwa \cs{frame}) wird dann verworfen, mit den
    % registrierten Ausnahmen. Entspricht genau dem Aufruf von \cs{mode*}
    % an derselben Stelle, an der Beamer diese Option selbst
    % auswertet.}{Activates the built-in \code{*} mode (see the model
    % defaults below): everything outside a unit environment (such as
    % \cs{frame}) is then dropped, with the registered exceptions.
    % Equivalent to calling \cs{mode*} at the same point where Beamer
    % itself evaluates this option.}
    ignorenonframetext .bool_gset:N =
      \g_osglecture_modes_ignorenonframetext_bool,
    ignorenonframetext .default:n = true,
    ignorenonframetext .initial:n = false,
    raw-skip-envs .clist_gset:N =
      \g_osglecture_modes_raw_skip_envs_clist,
    % \ldeen{Wird von einer Klasse gesetzt, die dieses Paket noch vor ihrem
    % eigenen \cs{LoadClass} lädt (osglecture tut das, damit der Modusgraph
    % rechtzeitig für modusabhängige Klassenoptionen abgeschlossen ist): Die
    % Auflösung des Blattmodus und das Patchen der Overlays warten dann auf
    % \cs{LectureModesActivateOverlays}, aufgerufen sobald die Basisklasse
    % vorhanden ist.}{A class that loads this package before its own
    % \cs{LoadClass} (osglecture does, so the mode graph is finalized in
    % time for class options keyed by mode) sets this: leaf-mode resolution
    % and overlay patching then wait for \cs{LectureModesActivateOverlays},
    % called once the base class is present.}
    deferred .bool_gset:N = \g_osglecture_modes_deferred_bool,
    deferred .initial:n = false,
  }
\ProcessKeyOptions [ osglecture_modes ]

\prg_new_protected_conditional:Npnn \osglecture_modes_if_mutable:n #1 { TF }
  {
    \bool_if:NTF \g_osglecture_modes_frozen_bool
      {
        \msg_error:nnn { osglecture_modes } { frozen } {#1}
        \prg_return_false:
      }
      { \prg_return_true: }
  }

\cs_new_protected:Npn \osglecture_modes_declare:n #1
  {
    \osglecture_modes_if_mutable:nTF {#1}
      {
        \seq_if_in:NnF \g_osglecture_modes_known_seq {#1}
          { \seq_gput_right:Nn \g_osglecture_modes_known_seq {#1} }
      }
      { }
  }

\cs_new_protected:Npn \osglecture_modes_declare_parents:nn #1#2
  {
    \osglecture_modes_if_mutable:nTF {#1}
      {
        \osglecture_modes_declare:n {#1}
        \clist_set:Nn \l_tmpa_clist {#2}
        \clist_remove_duplicates:N \l_tmpa_clist
        \prop_gput:NnV \g_osglecture_modes_parents_prop {#1} \l_tmpa_clist
      }
      { }
  }

\keys_define:nn { osglecture_modes / declare }
  {
    parent  .code:n = \osglecture_modes_declare_parents:Vn
      \l_osglecture_modes_decl_mode_tl {#1},
    parents .code:n = \osglecture_modes_declare_parents:Vn
      \l_osglecture_modes_decl_mode_tl {#1},
  }
\tl_new:N \l_osglecture_modes_decl_mode_tl
\cs_generate_variant:Nn \osglecture_modes_declare_parents:nn { Vn }
\cs_generate_variant:Nn \clist_map_inline:nn { Vn }

\NewDocumentCommand \DeclareLectureMode { O{} m }
  {
    \osglecture_modes_declare:n {#2}
    \tl_set:Nn \l_osglecture_modes_decl_mode_tl {#2}
    \keys_set:nn { osglecture_modes / declare } {#1}
  }

\NewDocumentCommand \DeclareLectureModeParents { m m }
  { \osglecture_modes_declare_parents:nn {#1} {#2} }

\NewDocumentCommand \DeclareLectureModeAlias { m m }
  {
    \osglecture_modes_if_mutable:nTF {#1}
      { \prop_gput:Nnn \g_osglecture_modes_alias_prop {#1} {#2} }
      { }
  }

\seq_new:N \l_osglecture_modes_alias_path_seq
\cs_new_protected:Npn \osglecture_modes_resolve:nN #1#2
  {
    \seq_clear:N \l_osglecture_modes_alias_path_seq
    \tl_set:Nn #2 {#1}
    \osglecture_modes_resolve_loop:N #2
  }
\cs_new_protected:Npn \osglecture_modes_resolve_loop:N #1
  {
    \prop_get:NVNT \g_osglecture_modes_alias_prop #1 \l_tmpa_tl
      {
        \seq_if_in:NVTF \l_osglecture_modes_alias_path_seq #1
          {
            \msg_error:nnV { osglecture_modes } { alias-cycle } #1
            \tl_clear:N #1
          }
          {
            \seq_put_right:NV \l_osglecture_modes_alias_path_seq #1
            \tl_set_eq:NN #1 \l_tmpa_tl
            \osglecture_modes_resolve_loop:N #1
          }
      }
  }

\cs_new_protected:Npn \osglecture_modes_activate:n #1
  {
    \osglecture_modes_if_mutable:nTF {#1}
      {
        \clist_map_inline:nn {#1}
          {
            \osglecture_modes_resolve:nN {##1} \l_tmpa_tl
            \seq_if_in:NVTF \g_osglecture_modes_known_seq \l_tmpa_tl
              {
                \seq_if_in:NVF \g_osglecture_modes_seed_seq \l_tmpa_tl
                  { \seq_gput_right:NV \g_osglecture_modes_seed_seq \l_tmpa_tl }
              }
              { \msg_error:nnV { osglecture_modes } { unknown-mode } \l_tmpa_tl }
          }
      }
      { }
  }

\NewDocumentCommand \AddLectureModes { m }
  { \osglecture_modes_activate:n {#1} }

% \ldeen{Kompatibilität: Das Setzen eines Modus ersetzt die vorherige
% Startmenge.}{Compatibility: setting a mode replaces the previous seed set.}
\NewDocumentCommand \SetLectureModes { m }
  {
    \osglecture_modes_if_mutable:nTF {#1}
      {
        \seq_gclear:N \g_osglecture_modes_seed_seq
        \osglecture_modes_activate:n {#1}
      }
      { }
  }

\NewExpandableDocumentCommand \LectureModes { }
  { \seq_use:Nn \g_osglecture_modes_seed_seq {,} }

\prop_new:N \l_osglecture_modes_visiting_prop
\prop_new:N \l_osglecture_modes_visited_prop
\bool_new:N \l_osglecture_modes_mark_active_bool
\cs_new_protected:Npn \osglecture_modes_visit:n #1
  {
    \prop_if_in:NnTF \l_osglecture_modes_visiting_prop {#1}
      { \msg_error:nnn { osglecture_modes } { graph-cycle } {#1} }
      {
        \prop_if_in:NnF \l_osglecture_modes_visited_prop {#1}
          {
            \prop_put:Nnn \l_osglecture_modes_visiting_prop {#1} { }
            \bool_if:NT \l_osglecture_modes_mark_active_bool
              { \prop_gput:Nnn \g_osglecture_modes_active_prop {#1} { } }
            \prop_get:NnNT \g_osglecture_modes_parents_prop {#1} \l_tmpa_tl
              {
                \clist_map_inline:Vn \l_tmpa_tl
                  {
                    \seq_if_in:NnTF \g_osglecture_modes_known_seq {##1}
                      { \osglecture_modes_visit:n {##1} }
                      {
                        \msg_error:nnnn { osglecture_modes } { unknown-parent }
                          {#1} {##1}
                      }
                  }
              }
            \bool_if:NT \l_osglecture_modes_mark_active_bool
              { \seq_gput_right:Nn \g_osglecture_modes_active_order_seq {#1} }
            \prop_remove:Nn \l_osglecture_modes_visiting_prop {#1}
            \prop_put:Nnn \l_osglecture_modes_visited_prop {#1} { }
          }
      }
  }

\cs_new_protected:Npn \osglecture_modes_validate_aliases:
  {
    \prop_map_inline:Nn \g_osglecture_modes_alias_prop
      {
        \osglecture_modes_resolve:nN {##1} \l_tmpa_tl
        \seq_if_in:NVF \g_osglecture_modes_known_seq \l_tmpa_tl
          {
            \msg_error:nnnn { osglecture_modes } { unknown-alias-target }
              {##1} {\l_tmpa_tl}
          }
      }
  }

\cs_new_protected:Npn \osglecture_modes_finalize:
  {
    \bool_if:NF \g_osglecture_modes_frozen_bool
      {
        \osglecture_modes_validate_aliases:
        \prop_gclear:N \g_osglecture_modes_active_prop
        \seq_gclear:N \g_osglecture_modes_active_order_seq
        \prop_clear:N \l_osglecture_modes_visiting_prop
        \prop_clear:N \l_osglecture_modes_visited_prop
        \bool_set_false:N \l_osglecture_modes_mark_active_bool
        \seq_map_inline:Nn \g_osglecture_modes_known_seq
          { \osglecture_modes_visit:n {##1} }
        \prop_clear:N \l_osglecture_modes_visiting_prop
        \prop_clear:N \l_osglecture_modes_visited_prop
        \bool_set_true:N \l_osglecture_modes_mark_active_bool
        \seq_map_inline:Nn \g_osglecture_modes_seed_seq
          { \osglecture_modes_visit:n {##1} }
        \bool_gset_true:N \g_osglecture_modes_frozen_bool
      }
  }
\NewDocumentCommand \FinalizeLectureModes { } { \osglecture_modes_finalize: }
\NewExpandableDocumentCommand \ActiveLectureModes { }
  { \seq_use:Nn \g_osglecture_modes_active_order_seq {,} }

\cs_generate_variant:Nn \osglecture_modes_visit:n { V }
% \ldeen{Aktiviert einen einzelnen, bereits deklarierten Modus zusätzlich
% zur aktuellen aktiven Menge -- anders als \cs{AddLectureModes} auch
% dann sicher, wenn der Graph schon abgeschlossen ist (finalisiert bei
% Bedarf zunächst, wie die erste Abfrage es ohnehin täte). \cs{osglecture_modes_visit:n}
% selbst ist dafür unverändert wiederverwendbar: Nach dem Abschluss bleibt
% \code{mark\_active\_bool} wahr und die Besucht-Buchführung bestehen,
% sodass ein erneuter Aufruf nur die noch nicht besuchten Vorfahren des
% neuen Modus markiert, bereits aktive Knoten aber unangetastet
% lässt.}{Activates a single, already-declared mode in addition to the
% current active set -- unlike \cs{AddLectureModes}, safe even once the
% graph is already finalized (finalizing first if needed, as the first
% query would anyway). \cs{osglecture_modes_visit:n} itself is directly
% reusable for this: after finalization, \code{mark\_active\_bool} stays
% true and the visited bookkeeping survives, so a further call only marks
% the new mode's not-yet-visited ancestors, leaving already-active nodes
% untouched.}
\cs_new_protected:Npn \osglecture_modes_activate_runtime:n #1
  {
    \osglecture_modes_resolve:nN {#1} \l_tmpa_tl
    \seq_if_in:NVTF \g_osglecture_modes_known_seq \l_tmpa_tl
      {
        \bool_if:NF \g_osglecture_modes_frozen_bool { \osglecture_modes_finalize: }
        \osglecture_modes_visit:V \l_tmpa_tl
      }
      { \msg_error:nnV { osglecture_modes } { unknown-mode } \l_tmpa_tl }
  }

% \ldeen{\cs{ModeValue} löst eine Spezifikation der Form
% \code{modus=wert|modus=wert|...} auf, indem \cs{ActiveLectureModes} von
% links nach rechts durchlaufen wird, mit \code{default} als letztem, stets
% vorhandenem Kandidaten -- dieselbe Reihenfolge und Priorität, die auch
% \cs{__semcat_resolve_presentation:nNNN} verwendet, hier aber expandierbar
% formuliert, damit das Ergebnis auch in \cs{dimen}-, \cs{skip}- und
% ähnliche Zuweisungen passt, nicht nur in Schlüssel-Wert-Schnittstellen.
% Eine Spezifikation ohne passenden Schlüssel und ohne \code{default}
% expandiert zu nichts.}{\cs{ModeValue} resolves a "mode=value|mode=value|..."
% spec by walking \cs{ActiveLectureModes} left to right, then trying
% "default" as the final, always-present candidate -- the same order and
% priority \cs{__semcat_resolve_presentation:nNNN} uses, just expressed
% expandably so the result drops into \cs{dimen}, \cs{skip}, \cs{count} and
% similar assignments, not only key-value interfaces. A spec without a
% matching key and without "default" expands to nothing.}
%
% \ldeen{Die Aufteilung erfolgt anhand von Trennzeichen, nicht durch echtes
% Schlüssel-Wert-Parsing, daher darf ein Wert weitere \code{=}-Zeichen oder
% Doppelpunkte enthalten (etwa eine URL), aber kein unmaskiertes \code{|};
% führende und folgende Leerzeichen um einen Schlüssel werden entfernt.
% \code{=} -- statt des sonst im Paket für eine \code{modus:rest}-Qualifikation
% genutzten Doppelpunkts -- ist hier der Schlüssel-Wert-Trenner, weil es von
% \cs{ExplSyntaxOn} unberührt bleibt, während ein Doppelpunkt in diesem
% Syntaxregime zum Buchstaben wird und daher nie auf den
% Kategorie-12-Doppelpunkt träfe, den ein Wert legitim enthalten könnte.
% Ein Segment ganz ohne \code{=} gilt als unqualifizierter Wert und steht für
% \code{default=} diesen Wert -- das stellt, unter \code{=}, dieselbe Kurzform
% wieder her, die die frühere, wegen Catcode-Problemen mit
% \cs{ExplSyntaxOn} aufgegebene \code{:}-basierte Spezifikationssyntax des
% Pakets kostenlos bot (dort war ein Segment ohne \code{modus:}-Präfix
% automatisch unqualifiziert).}{Splitting is delimiter-based, not
% key-value parsing, so a value may contain further "=" signs or colons
% (an URL, say) but not an unescaped "|"; leading and trailing spaces
% around a key are trimmed. "=" (rather than the ":" the rest of this
% package uses for a "mode:tail" overlay qualification) is the key/value
% separator here because it is untouched by \cs{ExplSyntaxOn}, unlike ":",
% which becomes a letter and so could never delimit on the category-12
% colon a value might legitimately contain. A segment carrying no "=" at
% all is an unqualified value and stands for "default=" that value --
% restoring, under "=", the convenience the package's former ":"-based
% spec syntax offered for free (there, an unqualified segment simply
% wasn't "mode:"-prefixed).}
\cs_new:Npn \osglecture_modes_value_seg:w #1#2#3#4|#5\q_stop
  {
    % \ldeen{\#1 = Kandidatenmodus ; \#2 = verbleibende Kandidaten ; \#3 =
    % Spezifikation-mit-"|" (für den nächsten Kandidaten aufbewahrt) ; \#4 =
    % dieses Segment ; \#5 = verbleibende Segmente}{\#1 = candidate mode ;
    % \#2 = remaining candidates ; \#3 = spec-with-"|" (kept for the next
    % candidate) ; \#4 = this segment ; \#5 = remaining segments}
    \osglecture_modes_value_kv_aux {#1} {#2} {#3} {#5} #4 \q_stop
  }

% \ldeen{\#5 ist hier der vollständige Rohtext dieses Segments, begrenzt
% nur durch das \cs{q_stop}, das \cs{osglecture_modes_value_seg:w} stets
% unmittelbar danach liefert -- anders als ein direktes Matching auf \code{=}
% endet dies also immer, ob das Segment eines enthält oder nicht. Ob es
% eines enthält, muss aber trotzdem getestet werden -- nur auf einer
% Wegwerfkopie, sodass die eigentliche Extraktion unten weiterhin einen
% Wert allein durch \cs{q_stop} begrenzen darf (er darf also weitere
% \code{=}-Zeichen enthalten): \code{=} \cs{q_no_value} an diese Kopie
% anzuhängen und dann auf das erste \code{=} zu matchen bedeutet, dass genau
% dann keines im Segment vorlag, wenn vor \cs{q_no_value} nichts mehr
% übrig ist, da ein bereits im Segment vorhandenes \code{=} stets mindestens
% das angehängte übrig lässt.}{\#5 here is this segment's full raw text,
% delimited only by the \cs{q_stop} \cs{osglecture_modes_value_seg:w}
% always supplies right after it -- so, unlike matching straight for a
% "=", this always terminates, whether or not the segment contains one.
% Whether it does still needs testing, but on a throwaway copy, so the
% real extraction below is free to keep delimiting a value by
% \cs{q_stop} alone (letting it contain further "=" signs): appending "="
% \cs{q_no_value} to that copy and matching the first "=" again means
% none was present in the segment exactly when nothing is left before
% \cs{q_no_value}, since a "=" already in the segment always leaves at
% least the appended one there.}
\cs_new:Npn \osglecture_modes_value_kv_aux #1#2#3#4#5\q_stop
  {
    \osglecture_modes_value_kv_test:w #5 = \q_no_value {#1} {#2} {#3} {#4} {#5}
  }

\cs_new:Npn \osglecture_modes_value_kv_test:w #1=#2\q_no_value #3#4#5#6#7
  {
    \tl_if_empty:nTF {#2}
      { \osglecture_modes_value_bare:nnnnn {#3} {#4} {#5} {#6} {#7} }
      { \osglecture_modes_value_kv_extract:w #7 \q_stop {#3} {#4} {#5} {#6} }
  }

\cs_new:Npn \osglecture_modes_value_kv_extract:w #1=#2\q_stop #3#4#5#6
  {
    % \ldeen{\#1 = Schlüssel dieses Segments ; \#2 = sein Wert (darf selbst
    % \code{=}-Zeichen enthalten, da er nur durch \cs{q_stop} begrenzt
    % wurde)}{\#1 = this segment's key ; \#2 = its value (may itself
    % contain "=" signs, since it was delimited only by \cs{q_stop})}
    \str_if_eq:eeTF { \tl_trim_spaces:n {#1} } {#3}
      { \tl_trim_spaces:n {#2} }
      { \osglecture_modes_value_continue:nnnn {#3} {#4} {#5} {#6} }
  }

% \ldeen{Ein unqualifiziertes (schlüsselloses) Segment gilt nur für den
% Kandidaten \code{default} als Treffer -- den stets vorhandenen letzten
% Ausweg -- und verhält sich damit genau wie ein explizites
% \code{default=...}-Segment; es kommt also nie einem Modus zuvor, der
% weiter vorn in \cs{ActiveLectureModes} passt.}{An unqualified (key-less)
% segment counts as a match only for the "default" candidate -- the
% always-present last resort -- so it behaves exactly like an explicit
% "default=..." entry and never pre-empts a mode that matches earlier in
% \cs{ActiveLectureModes}.}
\cs_new:Npn \osglecture_modes_value_bare:nnnnn #1#2#3#4#5
  {
    \str_if_eq:eeTF {#1} {default}
      { \tl_trim_spaces:n {#5} }
      { \osglecture_modes_value_continue:nnnn {#1} {#2} {#3} {#4} }
  }

\cs_new:Npn \osglecture_modes_value_continue:nnnn #1#2#3#4
  {
    \tl_if_empty:nTF {#4}
      { \osglecture_modes_value_next:nn {#2} {#3} }
      { \osglecture_modes_value_seg:w {#1} {#2} {#3} #4 \q_stop }
  }

\cs_new:Npn \osglecture_modes_value_next:nn #1#2
  {
    \tl_if_empty:nTF {#1}
      { }
      { \osglecture_modes_value_pop:w #1 \q_stop {#2} }
  }

% \ldeen{Entnimmt einen Kandidaten aus einer Kommaliste, die -- wie die
% Pipe-Liste oben -- stets ein synthetisches, angehängtes \code{,} trägt
% (wieder aus \cs{osglecture_modes_value_start:nn} bzw. bei jedem
% weiteren Pop von dieser Funktion selbst reproduziert), sodass auch
% \code{default}, der letzte Kandidat, keinen Sonderfall braucht.}{Pops one
% candidate off a comma list that -- like the pipe-list above -- always
% carries a synthetic trailing "," (again from
% \cs{osglecture_modes_value_start:nn}, or reproduced by this function's
% own output on every further pop), so "default", the last candidate,
% needs no special case either.}
\cs_new:Npn \osglecture_modes_value_pop:w #1,#2\q_stop #3
  { \osglecture_modes_value_seg:w {#1} {#2} {#3} #3 \q_stop }

\cs_new:Npn \osglecture_modes_value_priority: { \ActiveLectureModes , default }

\cs_new:Npn \osglecture_modes_value_start:nn #1#2
  { \osglecture_modes_value_pop:w #1 , \q_stop {#2|} }
\cs_generate_variant:Nn \osglecture_modes_value_start:nn { fn }

\NewExpandableDocumentCommand \ModeValue { m }
  { \osglecture_modes_value_start:fn { \osglecture_modes_value_priority: } {#1} }

\prg_new_conditional:Npnn \osglecture_modes_if_active:n #1 { T, F, TF }
  {
    \bool_if:NF \g_osglecture_modes_frozen_bool { \osglecture_modes_finalize: }
    \bool_set_false:N \l_tmpa_bool
    \tl_set:Nn \l_tmpa_tl {#1}
    % \ldeen{\code{||} ist der portable Oder-Operator; ein bloßes \code{|} ist
    % Backend-Overlaysyntax und hat in einem portablen Modusausdruck keine
    % Bedeutung.}{"||" is the portable disjunction operator; a bare "|"
    % is backend overlay syntax and has no meaning in a portable mode
    % expression.}
    \tl_replace_all:Nnn \l_tmpa_tl {||} {,}
    \clist_map_inline:Vn \l_tmpa_tl
      {
        \str_if_eq:nnTF {##1} {all}
          { \bool_set_true:N \l_tmpa_bool }
          {
            \osglecture_modes_resolve:nN {##1} \l_tmpa_tl
            \prop_if_in:NVTF \g_osglecture_modes_active_prop \l_tmpa_tl
              { \bool_set_true:N \l_tmpa_bool }
              {
                \seq_if_in:NVF \g_osglecture_modes_known_seq \l_tmpa_tl
                  { \msg_error:nnV { osglecture_modes } { unknown-mode } \l_tmpa_tl }
              }
          }
      }
    \bool_if:NTF \l_tmpa_bool
      { \prg_return_true: } { \prg_return_false: }
  }

\NewDocumentCommand \IfLectureModeTF { m +m +m }
  { \osglecture_modes_if_active:nTF {#1} {#2} {#3} }
\NewDocumentCommand \IfLectureModeT { m +m }
  { \osglecture_modes_if_active:nT {#1} {#2} }
\NewDocumentCommand \IfLectureModeF { m +m }
  { \osglecture_modes_if_active:nF {#1} {#2} }

\NewDocumentCommand \lecturemode { d<> +m }
  {
    \IfNoValueTF {#1}
      {#2}
      { \osglecture_modes_if_active:nT {#1} {#2} }
  }

% \ldeen{Die modusabhängige Befehlsauswahl beruht bewusst auf
% Implementationsbefehlen. Die gewählte Implementierung liest den
% ursprünglichen Argumentstrom selbst, sodass beliebige Befehlssignaturen
% möglich bleiben, ohne dass dieses Paket sie nachbilden müsste.}{Mode-aware
% command dispatch is deliberately based on implementation commands. The
% selected implementation reads the original argument stream, so arbitrary
% command signatures remain possible without duplicating them in this
% package.}
\seq_new:N  \g_osglecture_modes_behavior_classes_seq
\prop_new:N \g_osglecture_modes_mode_behavior_prop
\seq_new:N  \g_osglecture_modes_aware_commands_seq
\prop_new:N \g_osglecture_modes_command_requirements_prop
\prop_new:N \g_osglecture_modes_command_implementations_prop
\bool_new:N \g_osglecture_modes_commands_frozen_bool
\tl_new:N   \l_osglecture_modes_selected_behavior_tl
\cs_generate_variant:Nn \prop_gput:Nnn { NxV, Nxx }
\cs_generate_variant:Nn \prop_put:Nnn { NxV }
\cs_generate_variant:Nn \prop_get:NnN { NxN }
\cs_generate_variant:Nn \prop_get:NnNTF { NxNTF }

\msg_new:nnn { osglecture_modes } { commands-frozen }
  { The~mode-aware~command~registry~is~already~finalized. }
\msg_new:nnn { osglecture_modes } { unknown-behavior }
  { Unknown~mode~behavior~class~'#1'. }
\msg_new:nnn { osglecture_modes } { command-exists }
  {
    Mode-aware~command~'\token_to_str:N #1'~already~exists.
    ~Use~a~new~semantic~command~or~an~explicitly~saved~implementation.
  }
\msg_new:nnn { osglecture_modes } { unknown-aware-command }
  { Mode-aware~command~'\token_to_str:N #1'~has~not~been~declared. }
\msg_new:nnnn { osglecture_modes } { missing-implementation }
  {
    Mode-aware~command~'#1'~has~no~implementation
    ~for~behavior~'#2'.
  }
  { Register~one~before~finalizing~the~mode-aware~command~registry. }
\msg_new:nnnn { osglecture_modes } { undefined-implementation }
  {
    Implementation~'#2'~for~mode-aware~command~'#1'~is~undefined.
  }
  { Define~the~implementation~command~before~finalizing~the~registry. }
\msg_new:nnn { osglecture_modes } { no-active-behavior }
  { No~behavior~class~is~assigned~to~the~active~leaf~modes. }
\msg_new:nnnn { osglecture_modes } { conflicting-behaviors }
  { Active~leaf~modes~select~both~behavior~'#1'~and~'#2'. }
  { Assign~all~simultaneously~active~leaf~modes~to~one~behavior~class. }

\cs_new_protected:Npn \osglecture_modes_commands_if_mutable:T #1
  {
    \bool_if:NTF \g_osglecture_modes_commands_frozen_bool
      { \msg_error:nn { osglecture_modes } { commands-frozen } }
      {#1}
  }

\NewDocumentCommand \DeclareModeBehaviorClass { m }
  {
    \osglecture_modes_commands_if_mutable:T
      {
        \seq_if_in:NnF \g_osglecture_modes_behavior_classes_seq {#1}
          { \seq_gput_right:Nn \g_osglecture_modes_behavior_classes_seq {#1} }
      }
  }

\NewDocumentCommand \AssignLectureModeBehavior { m m }
  {
    \osglecture_modes_commands_if_mutable:T
      {
        \seq_if_in:NnTF \g_osglecture_modes_behavior_classes_seq {#2}
          {
            \clist_map_inline:nn {#1}
              {
                \osglecture_modes_resolve:nN {##1} \l_tmpa_tl
                \seq_if_in:NVTF \g_osglecture_modes_known_seq \l_tmpa_tl
                  {
                    \prop_gput:NVn
                      \g_osglecture_modes_mode_behavior_prop \l_tmpa_tl {#2}
                  }
                  {
                    \msg_error:nnV
                      { osglecture_modes } { unknown-mode } \l_tmpa_tl
                  }
              }
          }
          { \msg_error:nnn { osglecture_modes } { unknown-behavior } {#2} }
      }
  }

\cs_new_protected:Npn \osglecture_modes_use_aware_command:N #1
  {
    \osglecture_modes_finalize_commands:
    #1
  }

\NewDocumentCommand \DeclareModeAwareCommand { m m }
  {
    \osglecture_modes_commands_if_mutable:T
      {
        \cs_if_exist:NTF #1
          { \msg_error:nnN { osglecture_modes } { command-exists } #1 }
          {
            \seq_gput_right:Nx \g_osglecture_modes_aware_commands_seq
              { \cs_to_str:N #1 }
            \clist_set:Nn \l_tmpa_clist {#2}
            \clist_remove_duplicates:N \l_tmpa_clist
            \prop_gput:NxV \g_osglecture_modes_command_requirements_prop
              { \cs_to_str:N #1 } \l_tmpa_clist
            \cs_new_protected:Npn #1
              { \osglecture_modes_use_aware_command:N #1 }
          }
      }
  }

\NewDocumentCommand \DeclareModeAwareCommandImplementation { m m m }
  {
    \osglecture_modes_commands_if_mutable:T
      {
        \seq_if_in:NxTF \g_osglecture_modes_aware_commands_seq
          { \cs_to_str:N #1 }
          {
            \seq_if_in:NnTF \g_osglecture_modes_behavior_classes_seq {#2}
              {
                \prop_gput:Nxx
                  \g_osglecture_modes_command_implementations_prop
                  { \cs_to_str:N #1 / #2 } { \cs_to_str:N #3 }
              }
              {
                \msg_error:nnn
                  { osglecture_modes } { unknown-behavior } {#2}
              }
          }
          {
            \msg_error:nnN
              { osglecture_modes } { unknown-aware-command } #1
          }
      }
  }

\cs_new_protected:Npn \osglecture_modes_validate_aware_command:n #1
  {
    \prop_get:NnN \g_osglecture_modes_command_requirements_prop
      {#1} \l_tmpa_tl
    \clist_map_inline:Vn \l_tmpa_tl
      {
        \seq_if_in:NnTF \g_osglecture_modes_behavior_classes_seq {##1}
          {
            \prop_get:NnNTF \g_osglecture_modes_command_implementations_prop
              {#1 / ##1} \l_tmpb_tl
              {
                \cs_if_exist:cF {\l_tmpb_tl}
                  {
                    \msg_error:nnxx
                      { osglecture_modes } { undefined-implementation }
                      {#1} {\l_tmpb_tl}
                  }
              }
              {
                \msg_error:nnxx
                  { osglecture_modes } { missing-implementation }
                  {#1} {##1}
              }
          }
          { \msg_error:nnn { osglecture_modes } { unknown-behavior } {##1} }
      }
  }

\cs_new_protected:Npn \osglecture_modes_select_behavior:
  {
    \tl_clear:N \l_osglecture_modes_selected_behavior_tl
    \seq_map_inline:Nn \g_osglecture_modes_seed_seq
      {
        \prop_get:NnNT \g_osglecture_modes_mode_behavior_prop
          {##1} \l_tmpa_tl
          {
            \tl_if_empty:NTF \l_osglecture_modes_selected_behavior_tl
              {
                \tl_set_eq:NN
                  \l_osglecture_modes_selected_behavior_tl \l_tmpa_tl
              }
              {
                \str_if_eq:VVF
                  \l_osglecture_modes_selected_behavior_tl \l_tmpa_tl
                  {
                    \msg_error:nnVV
                      { osglecture_modes } { conflicting-behaviors }
                      \l_osglecture_modes_selected_behavior_tl \l_tmpa_tl
                  }
              }
          }
      }
    \tl_if_empty:NT \l_osglecture_modes_selected_behavior_tl
      { \msg_error:nn { osglecture_modes } { no-active-behavior } }
  }

\cs_new_protected:Npn \osglecture_modes_bind_aware_command:n #1
  {
    \prop_get:NxNTF \g_osglecture_modes_command_implementations_prop
      {#1 / \l_osglecture_modes_selected_behavior_tl} \l_tmpa_tl
      { \cs_gset_eq:cc {#1} {\l_tmpa_tl} }
      {
        \msg_error:nnxx { osglecture_modes } { missing-implementation }
          {#1} {\l_osglecture_modes_selected_behavior_tl}
      }
  }

\cs_new_protected:Npn \osglecture_modes_finalize_commands:
  {
    \bool_if:NF \g_osglecture_modes_commands_frozen_bool
      {
        \osglecture_modes_finalize:
        \seq_map_inline:Nn \g_osglecture_modes_aware_commands_seq
          { \osglecture_modes_validate_aware_command:n {##1} }
        \seq_if_empty:NF \g_osglecture_modes_aware_commands_seq
          {
            \osglecture_modes_select_behavior:
            \seq_map_inline:Nn \g_osglecture_modes_aware_commands_seq
              { \osglecture_modes_bind_aware_command:n {##1} }
          }
        \bool_gset_true:N \g_osglecture_modes_commands_frozen_bool
      }
  }

\NewDocumentCommand \FinalizeModeAwareCommands { }
  { \osglecture_modes_finalize_commands: }

% \ldeen{Portable Modusspezifikationen werden hier ausgewertet. Echte
% Overlayspezifikationen nutzen ein natives \cs{alt}, das entweder den
% ursprünglichen Befehl oder einen signaturbewussten Schlucker wählt,
% bevor die restlichen Argumente gelesen werden.}{Portable mode
% specifications are evaluated here. Genuine overlay specifications use a
% native \cs{alt}, which selects either the original command or a
% signature-aware gobbler before the remaining arguments are read.}
\tl_new:N \l_osglecture_modes_overlay_default_tl
\tl_new:N \l_osglecture_modes_overlay_fallback_tl
\tl_new:N \l_osglecture_modes_overlay_arguments_tl
\prop_new:N \g_osglecture_modes_overlay_default_prop
\prop_new:N \g_osglecture_modes_overlay_fallback_prop
\msg_new:nnn { osglecture_modes } { overlay-fallback-error }
  { Overlay~specification~'#1'~requires~a~native~overlay~backend. }
\keys_define:nn { osglecture_modes / overlay-aware }
  {
    arguments .tl_set:N = \l_osglecture_modes_overlay_arguments_tl,
    default-mode .tl_set:N = \l_osglecture_modes_overlay_default_tl,
    fallback .choices:nn = { flatten, omit, error }
      { \tl_set_eq:NN \l_osglecture_modes_overlay_fallback_tl \l_keys_choice_tl },
    fallback .initial:n = flatten,
  }

\cs_new_protected:Npn \osglecture_modes_overlay_original:n #1
  { \use:c { osglecture_modes_overlay_original_#1 } }

\cs_new_protected:Npn \osglecture_modes_overlay_gobble:n #1
  { \use:c { osglecture_modes_overlay_gobble_#1 } }

\cs_new_protected:Npn \osglecture_modes_overlay_fallback:nn #1#2
  {
    \prop_get:NnN \g_osglecture_modes_overlay_fallback_prop {#1} \l_tmpa_tl
    \str_case:Vn \l_tmpa_tl
      {
        {flatten}{ \osglecture_modes_overlay_original:n {#1} }
        {omit}{ \osglecture_modes_overlay_gobble:n {#1} }
        {error}{
          \msg_error:nnn
            { osglecture_modes } { overlay-fallback-error } {#2}
          \osglecture_modes_overlay_gobble:n {#1}
        }
      }
  }

% \ldeen{Einmalig erfasst, bevor ein portabler \cs{alt}-Fallback weiter
% unten überhaupt existieren und fälschlich für ein natives Urteil
% gehalten werden könnte: \cs{alt} dient dieser Funktion zugleich als
% eigener Weg, ein echtes Overlaybackend zu fragen, ob eine beliebige
% (nicht-portable) Spezifikation aktiv ist -- das kann nur ein echtes
% Backend beantworten; ein rein portables \cs{alt} versteht benannte
% Modi, keine beliebigen Spezifikationen.}{Captured once, before any
% portable-only \cs{alt} fallback further below could exist and be
% mistaken for a native judge: \cs{alt} doubles as this function's own
% way of asking a real overlay backend whether an arbitrary (non-portable)
% spec is active, which only a genuine backend can answer -- a
% portable-only \cs{alt} understands named modes, not arbitrary specs.}
\bool_new:N \g_osglecture_modes_native_alt_bool
\bool_set:Nn \g_osglecture_modes_native_alt_bool { \cs_if_exist_p:N \alt }

\cs_new_protected:Npn \osglecture_modes_overlay_alt:nnn #1#2#3
  { \alt<#1>{#2}{#3} }
\cs_generate_variant:Nn \osglecture_modes_overlay_alt:nnn { Vnn }

\cs_new_protected:Npn \osglecture_modes_overlay_apply:nn #1#2
  {
    \osglecture_modes_classify:n {#2}
    \str_case:Vn \l_osglecture_modes_action_tl
      {
        { show } { \osglecture_modes_overlay_original:n {#1} }
        { hide } { \osglecture_modes_overlay_gobble:n {#1} }
        { native }
          {
            \bool_if:NTF \g_osglecture_modes_native_alt_bool
              {
                \osglecture_modes_overlay_alt:Vnn
                  \l_osglecture_modes_native_tl
                  { \osglecture_modes_overlay_original:n {#1} }
                  { \osglecture_modes_overlay_gobble:n {#1} }
              }
              { \osglecture_modes_overlay_fallback:nn {#1} {#2} }
          }
      }
  }

\NewDocumentCommand \MakeOverlayAwareCommand { O{} m }
  {
    \tl_clear:N \l_osglecture_modes_overlay_default_tl
    \tl_clear:N \l_osglecture_modes_overlay_arguments_tl
    \tl_set:Nn \l_osglecture_modes_overlay_fallback_tl {flatten}
    \keys_set:nn { osglecture_modes / overlay-aware } {#1}
    \cs_if_exist:NTF #2
      {
        \cs_if_exist:cTF
          { osglecture_modes_overlay_original_\cs_to_str:N #2 }
          {
            \exp_args:Nc \RenewCommandCopy
              { osglecture_modes_overlay_original_\cs_to_str:N #2 } #2
          }
          {
            \exp_args:Nc \NewCommandCopy
              { osglecture_modes_overlay_original_\cs_to_str:N #2 } #2
          }
        \exp_args:NcV \DeclareDocumentCommand
          { osglecture_modes_overlay_gobble_\cs_to_str:N #2 }
          \l_osglecture_modes_overlay_arguments_tl
          { }
        \prop_put:NxV \g_osglecture_modes_overlay_default_prop
          { \cs_to_str:N #2 } \l_osglecture_modes_overlay_default_tl
        \prop_put:NxV \g_osglecture_modes_overlay_fallback_prop
          { \cs_to_str:N #2 } \l_osglecture_modes_overlay_fallback_tl
        \RenewDocumentCommand #2 { d<> }
          {
            \IfNoValueTF {##1}
              {
                \prop_get:NxN \g_osglecture_modes_overlay_default_prop
                  { \cs_to_str:N #2 } \l_tmpa_tl
                \tl_if_empty:NTF \l_tmpa_tl
                  {
                    \osglecture_modes_overlay_original:n
                      { \cs_to_str:N #2 }
                  }
                  {
                    \exp_args:Nx \osglecture_modes_overlay_apply:nV
                      { \cs_to_str:N #2 }
                      \l_tmpa_tl
                  }
              }
              {
                \exp_args:Nx \osglecture_modes_overlay_apply:nn
                  { \cs_to_str:N #2 } {##1}
              }
          }
      }
      { \msg_error:nnN { kernel } { command-not-defined } #2 }
  }
\cs_generate_variant:Nn \osglecture_modes_overlay_apply:nn { nV }

% \ldeen{Jede an \cs{only}, \cs{mode}, die overlayfähigen Befehle oder die
% gepatchten Backend-Decoder übergebene Winkelspezifikation wird in eine
% von drei Arten einsortiert. Der Klassifizierer merkt sich anzeigen,
% verbergen oder nativ, und im nativen Fall zusätzlich die an das Backend
% weiterzureichende Spezifikation.
% \begin{itemize}
%   \item \emph{portabel} -- eine mit \code{||} verknüpfte Oder-Liste
%     registrierter Modi (oder \code{all}). Wird vollständig hier
%     ausgewertet; benötigt kein Overlaybackend.
%   \item \emph{modusqualifiziert} -- nicht portabel, aber jedes mit
%     \code{|} getrennte Segment ist ein registrierter Modus, ein
%     \code{registriertermodus:rest} oder ein rein numerisches Overlayatom.
%     Die portablen Anteile werden hier ausgewertet und die überlebenden
%     nativen Reste zu einer einzigen Backendspezifikation
%     zusammengesetzt.
%   \item \emph{nativ} -- alles Übrige: Backend-Modusnamen,
%     \code{|}-Modusskopierung über unbekannte Modi, Actionspezifikationen,
%     \dots{} Wird unverändert weitergereicht.
% \end{itemize}
% \code{|} und \code{,} innerhalb von \code{<...>} sind stets Backendsyntax.
% \code{||} ist der einzige portable Oder-Operator, daher muss eine
% Spezifikation, die portable Modi mit nativen Overlayatomen mischt, dies
% über \code{||} (oder die Form \code{modus:rest}) kenntlich machen.}{Every
% angle-bracket specification handed to \cs{only}, \cs{mode}, the
% overlay-aware commands, or the patched backend decoders is sorted into
% one of three kinds. The classifier records show, hide, or native, and
% for a native verdict the specification to forward to the backend.
% \begin{itemize}
%   \item \emph{portable} -- a "||"-separated disjunction of registered
%     modes (or "all"). Evaluated entirely here; needs no overlay
%     backend.
%   \item \emph{qualified} -- not portable, but every "|"-separated
%     segment is a registered mode, "registeredmode:tail", or a purely
%     numeric overlay atom. The portable parts are evaluated here and the
%     surviving native tails are recombined into one backend spec.
%   \item \emph{native} -- anything else: backend mode names, "|" mode
%     scoping over unknown modes, action specs, \dots{} Forwarded
%     verbatim.
% \end{itemize}
% "|" and "," inside \code{<...>} are always backend syntax. "||" is the
% only portable disjunction operator, so a spec that mixes portable modes
% with native overlay atoms must use "||" (or the "mode:tail" form) to
% say so.}
\tl_new:N    \l_osglecture_modes_action_tl
\tl_new:N    \l_osglecture_modes_native_tl
\bool_new:N  \l_osglecture_modes_spec_ok_bool
\bool_new:N  \l_osglecture_modes_spec_qual_bool
\seq_new:N   \l_osglecture_modes_spec_atoms_seq
\seq_new:N   \l_osglecture_modes_spec_segs_seq
\seq_new:N   \l_osglecture_modes_spec_colon_seq
\clist_new:N \l_osglecture_modes_spec_parts_clist
\tl_new:N    \l_osglecture_modes_spec_mode_tl
\tl_new:N    \l_osglecture_modes_spec_rest_tl
\tl_new:N    \l_osglecture_modes_spec_resolved_tl

\cs_generate_variant:Nn \osglecture_modes_resolve:nN { VN }
\prg_generate_conditional_variant:Nnn \osglecture_modes_if_active:n
  { V } { T, F, TF }
% \ldeen{\code{:} ist in diesem Syntaxregime ein Buchstabe, ein Literal
% trifft daher nie auf die Kategorie-12-Doppelpunkte, die uns aus dem
% Dokument erreichen; stattdessen anhand der Doppelpunkt-Stringkonstante
% aufteilen.}{":" is a letter in this syntax regime, so a literal one
% never matches the category-12 colons that reach us from the document;
% split on the colon string constant instead.}
\cs_generate_variant:Nn \seq_set_split:Nnn { NVn }
\cs_generate_variant:Nn \seq_use:Nn { NV }

\prg_new_protected_conditional:Npnn \osglecture_modes_spec_if_portable:n #1
  { TF }
  {
    \seq_set_split:Nnn \l_osglecture_modes_spec_atoms_seq { || } {#1}
    \bool_set_true:N \l_osglecture_modes_spec_ok_bool
    \seq_map_inline:Nn \l_osglecture_modes_spec_atoms_seq
      {
        \tl_set:Nn \l_osglecture_modes_spec_mode_tl {##1}
        \tl_trim_spaces:N \l_osglecture_modes_spec_mode_tl
        \str_if_eq:VnF \l_osglecture_modes_spec_mode_tl { all }
          {
            \osglecture_modes_resolve:VN \l_osglecture_modes_spec_mode_tl
              \l_osglecture_modes_spec_resolved_tl
            \seq_if_in:NVF \g_osglecture_modes_known_seq
              \l_osglecture_modes_spec_resolved_tl
              { \bool_set_false:N \l_osglecture_modes_spec_ok_bool }
          }
      }
    \bool_if:NTF \l_osglecture_modes_spec_ok_bool
      { \prg_return_true: } { \prg_return_false: }
  }

% \ldeen{Ein Segment hat \emph{native Form}, wenn es nur aus
% Overlay-Zahlzeichen besteht -- es ist dann ein modusunabhängiges
% Overlayatom statt eines Backend-Modusnamens.}{A segment is "native
% shape" when it holds only overlay-number characters, i.e. it is a
% mode-independent overlay atom rather than a backend mode name.}
\prg_new_protected_conditional:Npnn \osglecture_modes_spec_if_native_shape:n
  #1 { TF }
  {
    \bool_set_true:N \l_osglecture_modes_spec_ok_bool
    \str_map_inline:nn {#1}
      {
        \str_if_in:nnF { 0123456789+-.(),~ } {##1}
          {
            \bool_set_false:N \l_osglecture_modes_spec_ok_bool
            \str_map_break:
          }
      }
    \bool_if:NTF \l_osglecture_modes_spec_ok_bool
      { \prg_return_true: } { \prg_return_false: }
  }

\cs_new_protected:Npn \osglecture_modes_spec_segment_bare:n #1
  {
    \tl_set:Nn \l_osglecture_modes_spec_mode_tl {#1}
    \osglecture_modes_resolve:VN \l_osglecture_modes_spec_mode_tl
      \l_osglecture_modes_spec_resolved_tl
    \seq_if_in:NVTF \g_osglecture_modes_known_seq
      \l_osglecture_modes_spec_resolved_tl
      {
        % \ldeen{Ein bloßer registrierter Modus bedeutet \emph{in diesem
        % Modus, auf jeder Folie}.}{A bare registered mode means "in
        % this mode, on every slide".}
        \osglecture_modes_if_active:VT \l_osglecture_modes_spec_resolved_tl
          { \clist_put_right:Nn \l_osglecture_modes_spec_parts_clist { all } }
      }
      {
        \osglecture_modes_spec_if_native_shape:nTF {#1}
          { \clist_put_right:Nn \l_osglecture_modes_spec_parts_clist {#1} }
          { \bool_set_false:N \l_osglecture_modes_spec_qual_bool }
      }
  }

\cs_new_protected:Npn \osglecture_modes_spec_segment_colon:n #1
  {
    \seq_set_split:NVn \l_osglecture_modes_spec_colon_seq \c_colon_str {#1}
    \seq_pop_left:NN \l_osglecture_modes_spec_colon_seq
      \l_osglecture_modes_spec_mode_tl
    \tl_set:Nx \l_osglecture_modes_spec_rest_tl
      { \seq_use:NV \l_osglecture_modes_spec_colon_seq \c_colon_str }
    \tl_trim_spaces:N \l_osglecture_modes_spec_mode_tl
    \tl_trim_spaces:N \l_osglecture_modes_spec_rest_tl
    \osglecture_modes_resolve:VN \l_osglecture_modes_spec_mode_tl
      \l_osglecture_modes_spec_resolved_tl
    \seq_if_in:NVTF \g_osglecture_modes_known_seq
      \l_osglecture_modes_spec_resolved_tl
      {
        \osglecture_modes_if_active:VT \l_osglecture_modes_spec_resolved_tl
          {
            \clist_put_right:NV \l_osglecture_modes_spec_parts_clist
              \l_osglecture_modes_spec_rest_tl
          }
      }
      { \bool_set_false:N \l_osglecture_modes_spec_qual_bool }
  }

\cs_new_protected:Npn \osglecture_modes_spec_segment:n #1
  {
    \bool_if:NT \l_osglecture_modes_spec_qual_bool
      {
        \tl_set:Nn \l_osglecture_modes_spec_mode_tl {#1}
        \tl_trim_spaces:N \l_osglecture_modes_spec_mode_tl
        \tl_if_empty:NF \l_osglecture_modes_spec_mode_tl
          {
            \exp_args:NV \str_if_in:nnTF \l_osglecture_modes_spec_mode_tl { : }
              {
                \exp_args:NV \osglecture_modes_spec_segment_colon:n
                  \l_osglecture_modes_spec_mode_tl
              }
              {
                \exp_args:NV \osglecture_modes_spec_segment_bare:n
                  \l_osglecture_modes_spec_mode_tl
              }
          }
      }
  }

\cs_new_protected:Npn \osglecture_modes_classify:n #1
  {
    \tl_clear:N \l_osglecture_modes_native_tl
    \osglecture_modes_spec_if_portable:nTF {#1}
      {
        \osglecture_modes_if_active:nTF {#1}
          { \tl_set:Nn \l_osglecture_modes_action_tl { show } }
          { \tl_set:Nn \l_osglecture_modes_action_tl { hide } }
      }
      {
        \bool_set_true:N \l_osglecture_modes_spec_qual_bool
        \clist_clear:N \l_osglecture_modes_spec_parts_clist
        \seq_set_split:Nnn \l_osglecture_modes_spec_segs_seq { | } {#1}
        \seq_map_inline:Nn \l_osglecture_modes_spec_segs_seq
          { \osglecture_modes_spec_segment:n {##1} }
        \bool_if:NTF \l_osglecture_modes_spec_qual_bool
          {
            \clist_if_empty:NTF \l_osglecture_modes_spec_parts_clist
              { \tl_set:Nn \l_osglecture_modes_action_tl { hide } }
              {
                \tl_set:Nn \l_osglecture_modes_action_tl { native }
                \tl_set:Nx \l_osglecture_modes_native_tl
                  { \clist_use:Nn \l_osglecture_modes_spec_parts_clist { , } }
              }
          }
          {
            \tl_set:Nn \l_osglecture_modes_action_tl { native }
            \tl_set:Nn \l_osglecture_modes_native_tl {#1}
          }
      }
  }

% \ldeen{Inhalt wird angezeigt, außer ein portabler Ausdruck ist vorhanden
% und wertet zu falsch aus; eine native oder unerkannte Spezifikation
% degradiert dort zu \emph{anzeigen}, wo kein Overlaybackend sie erfüllen
% kann.}{Content shows unless a portable expression is present and
% evaluates false; a native or unrecognised specification degrades to
% "show" where no overlay backend can honour it.}
\prg_new_protected_conditional:Npnn \osglecture_modes_if_reveal:n #1
  { T, F, TF }
  {
    \osglecture_modes_classify:n {#1}
    \str_if_eq:VnTF \l_osglecture_modes_action_tl { hide }
      { \prg_return_false: } { \prg_return_true: }
  }

% \ldeen{Die Vorgaben beschreiben Semantik, keine Backend-Wahl.}{Defaults
% describe semantics, not a backend choice.}
\DeclareLectureMode{presentation}
\DeclareLectureMode[parents=presentation]{beamer}
\DeclareLectureMode[parents=presentation]{projector}
\DeclareLectureMode[parents={presentation,print}]{handout}
\DeclareLectureMode[parents=presentation]{trans}
\DeclareLectureMode{longform}
\DeclareLectureMode[parents={longform,print}]{article}
\DeclareLectureMode[parents=article]{script}
\DeclareLectureMode{print}
\DeclareLectureMode[parents=longform]{web}
\DeclareLectureModeAlias{slides}{projector}

% \ldeen{Der Modus \code{*} ist ein gewöhnlicher, aber vorinstallierter
% Modus: Ist er aktiv, wird alles außerhalb einer registrierten
% Unit-Umgebung (per Vorgabe \cs{frame}, falls vorhanden) verworfen, mit
% den unten registrierten Ausnahmen. Aktiviert wird er entweder über die
% Paketoption \code{ignorenonframetext} (Vorgabe: \code{true}, s.\,o.)
% oder -- wie jeder andere Filtermodus auch -- direkt über \cs{mode*}
% beziehungsweise \cs{mode<*>}; \cs{mode*} ist dabei nur eine wegen
% Beamer-Kompatibilität beibehaltene Syntaxvariante für
% \cs{mode<*>}, kein eigener Mechanismus.}{Mode \code{*} is an ordinary,
% but pre-installed mode: while active, everything outside a registered
% unit environment (by default \cs{frame}, if present) is dropped, with
% the exceptions registered below. It is activated either via the
% \code{ignorenonframetext} package option (default: \code{true}, see
% above) or, like any other filter mode, directly via \cs{mode*} or
% \cs{mode<*>}; \cs{mode*} is only a syntax variant kept for Beamer
% compatibility with \cs{mode<*>}, not a separate mechanism.}
\DeclareLectureMode{*}

% \ldeen{Modi, die über die Schalterform \cs{mode<name>} (ohne
% Inhaltsargument) statt über \cs{lecturemode}/\cs{IfLectureModeTF}
% eingesetzt werden, heißen hier \emph{Filtermodi}: Ihre Aktivierung
% schaltet nicht Inhalt an einer Stelle sichtbar, sondern engagiert --
% global, ab dieser Stelle -- das Verwerfen alles Unregistrierten
% außerhalb der registrierten Unit-Umgebungen. Sie werden wie jeder
% andere Modus mit \cs{DeclareLectureMode} deklariert; nichts an ihrer
% Deklaration unterscheidet sie von inhaltlichen Modi wie \code{script}
% oder \code{presentation}. Ausnahmen (welche Befehle/Umgebungen dabei
% erhalten bleiben) werden separat über \cs{DeclareLectureKeepCommand},
% \cs{DeclareLectureKeepLabelCommand} und \cs{DeclareLectureKeepEnvironment}
% registriert, optional an einen bestimmten Filtermodus gebunden
% (\code{mode=\meta{name}}); ohne diesen Schlüssel gilt die Ausnahme als
% Baseline, unabhängig davon, welcher Filtermodus (falls überhaupt einer)
% gerade aktiv ist.}{Modes used via the switch form \cs{mode<name>}
% (without a content argument) rather than via
% \cs{lecturemode}/\cs{IfLectureModeTF} are called \emph{filter modes}
% here: activating one does not toggle content visible at one spot, but
% engages -- globally, from that point on -- dropping everything
% unregistered outside the registered unit environments. They are
% declared like any other mode, with \cs{DeclareLectureMode}; nothing
% about their declaration distinguishes them from content modes such as
% \code{script} or \code{presentation}. Exceptions (which
% commands/environments survive) are registered separately via
% \cs{DeclareLectureKeepCommand}, \cs{DeclareLectureKeepLabelCommand},
% and \cs{DeclareLectureKeepEnvironment}, optionally bound to a specific
% filter mode (\code{mode=\meta{name}}); without that key, an exception
% counts as baseline, in effect no matter which filter mode, if any, is
% currently active.}
\bool_new:N  \g_osglecture_modes_scanner_active_bool
\bool_new:N  \g_osglecture_modes_after_structure_bool
\bool_new:N  \g_osglecture_modes_scan_engaged_bool
\bool_new:N  \g_osglecture_modes_scan_extended_bool
\bool_new:N  \g_osglecture_modes_native_mode_exists_bool
\clist_new:N \g_osglecture_modes_keep_envs_clist
\clist_new:N \g_osglecture_modes_hooked_envs_clist
\prop_new:N  \g_osglecture_modes_dispatch_prop
\prop_new:N  \g_osglecture_modes_dispatch_mode_prop
\prop_new:N  \g_osglecture_modes_keep_envs_mode_prop
\tl_new:N    \l_osglecture_modes_keep_mode_tl

\msg_new:nnn { osglecture_modes } { unknown-environment }
  { Environment~'#1'~is~not~defined. }

\keys_define:nn { osglecture_modes / keep }
  {
    mode .tl_set:N = \l_osglecture_modes_keep_mode_tl,
  }

\sys_if_engine_luatex:T
  { \directlua{require("osglecture-modes-ltxtalk.lua")} }

\cs_new_protected:Npn \osglecture_modes_lua_raw_skip:n #1
  {
    \sys_if_engine_luatex:T
      {
        \directlua{
          osglecture_modes_ltxtalk.start_raw_skip("\luaescapestring{#1}")
        }
      }
  }

% \ldeen{Engagiert den Scanner nur, wenn ihn zuvor \cs{osglecture_modes_switch:n}
% tatsächlich eingeschaltet hat (weil kein natives \cs{mode*} zur
% Verfügung stand oder eigene Ausnahmen über die native Positivliste
% hinausgehen), und nur außerhalb einer gerade betretenen
% Unit-Umgebung.}{Only engages the scanner if
% \cs{osglecture_modes_switch:n} actually turned it on (because no
% native \cs{mode*} was available, or because custom exceptions go
% beyond the native allow-list), and only outside a unit environment
% just entered.}
\cs_new_protected:Npn \osglecture_modes_start_scan:
  {
    \bool_if:NT \g_osglecture_modes_scan_engaged_bool
      {
        \bool_if:NT \g_osglecture_modes_scanner_active_bool
          { \osglecture_modes_scan: }
      }
  }

\cs_new_protected:Npn \osglecture_modes_scan:
  {
    \peek_remove_spaces:n
      {
        \peek_meaning:NTF \begin
          { \osglecture_modes_scan_begin: }
          {
            \peek_meaning:NTF \end
              { \osglecture_modes_scan_end: }
              { \osglecture_modes_scan_token: }
          }
      }
  }

\cs_new_protected:Npn \osglecture_modes_scan_begin:
  {
    \peek_meaning_remove:NT \begin
    \osglecture_modes_scan_begin_aux:n
  }

\cs_new_protected:Npn \osglecture_modes_scan_begin_aux:n #1
  {
    \clist_if_in:NnTF \g_osglecture_modes_keep_envs_clist {#1}
      {
        \prop_get:NnNTF \g_osglecture_modes_keep_envs_mode_prop {#1} \l_tmpa_tl
          {
            \osglecture_modes_if_active:VTF \l_tmpa_tl
              { \begin{#1} }
              {
                \osglecture_modes_skip_environment:n {#1}
                \osglecture_modes_scan:
              }
          }
          { \begin{#1} }
      }
      {
        \clist_if_in:NnTF \g_osglecture_modes_raw_skip_envs_clist {#1}
          {
            \osglecture_modes_lua_raw_skip:n {#1}
            \osglecture_modes_scan:
          }
          {
            \osglecture_modes_skip_environment:n {#1}
            \osglecture_modes_scan:
          }
      }
  }

\cs_new_protected:Npn \osglecture_modes_scan_end:
  {
    \peek_meaning_remove:NT \end
    \osglecture_modes_scan_end_aux:n
  }

\cs_new_protected:Npn \osglecture_modes_scan_end_aux:n #1
  {
    \str_if_eq:nnTF {#1} {document}
      {
        \end{document}
      }
      {
        % \ldeen{fremdes \cs{end}\code{\{...\}} außerhalb von Unit-Umgebungen
        % ignorieren}{ignore a foreign \cs{end}\code{\{...\}} outside unit
        % environments}
        \osglecture_modes_scan:
      }
  }

\cs_new_protected:Npn \osglecture_modes_skip_environment:n #1
  {
    \cs_set_protected:Npn \osglecture_modes_skip:w ##1 \end ##2
      {
        \str_if_eq:nnF {#1} {##2}
          { \osglecture_modes_skip:w }
      }
    \osglecture_modes_skip:w
  }

\cs_new_protected:Npn \osglecture_modes_ensure_env_hook:n #1
  {
    \cs_if_exist:cTF {#1}
      {
        \clist_if_in:NnF \g_osglecture_modes_hooked_envs_clist {#1}
          {
            \clist_gput_right:Nn \g_osglecture_modes_hooked_envs_clist {#1}
            \AddToHook{env/#1/after}{ \osglecture_modes_start_scan: }
          }
      }
      { \msg_error:nnn { osglecture_modes } { unknown-environment } {#1} }
  }

\cs_new_protected:Npn \osglecture_modes_declare_keep_environment:nn #1#2
  {
    \osglecture_modes_ensure_env_hook:n {#1}
    \clist_if_in:NnF \g_osglecture_modes_keep_envs_clist {#1}
      { \clist_gput_right:Nn \g_osglecture_modes_keep_envs_clist {#1} }
    \tl_if_empty:nF {#2}
      { \prop_gput:Nnn \g_osglecture_modes_keep_envs_mode_prop {#1} {#2} }
  }
\cs_generate_variant:Nn \osglecture_modes_declare_keep_environment:nn { nV }

\NewDocumentCommand \DeclareLectureKeepEnvironment { O{} m }
  {
    \tl_clear:N \l_osglecture_modes_keep_mode_tl
    \keys_set:nn { osglecture_modes / keep } {#1}
    \bool_gset_true:N \g_osglecture_modes_scan_extended_bool
    \clist_map_inline:nn {#2}
      {
        \osglecture_modes_declare_keep_environment:nV
          {##1} \l_osglecture_modes_keep_mode_tl
      }
  }

\cs_if_exist:cT { frame }
  { \osglecture_modes_declare_keep_environment:nn { frame } { } }
\cs_if_exist:cT { frame* }
  { \osglecture_modes_declare_keep_environment:nn { frame* } { } }

% \ldeen{Eine ohne \code{mode=} deklarierte Ausnahme gilt als Baseline:
% Sie erweitert die native Positivliste immer, deshalb erzwingt sie hier
% wie eine modusgebundene Ausnahme den eigenen Scanner (siehe
% \cs{osglecture_modes_switch:n} weiter unten) -- die intern registrierten
% Basisausnahmen (Gliederungsbefehle, \cs{mode}, \cs{lecturemode},
% \cs{maketitle}, \cs{frame}) setzen diesen Schalter dagegen bewusst
% nicht: Sie spiegeln genau das, was Beamers eigenes \cs{mode*} ohnehin
% kennt.}{An exception declared without \code{mode=} counts as baseline:
% it always extends the native allow-list, so it forces the package's
% own scanner here just like a mode-bound exception (see
% \cs{osglecture_modes_switch:n} further below) -- the internally
% registered baseline exceptions (structural commands, \cs{mode},
% \cs{lecturemode}, \cs{maketitle}, \cs{frame}) deliberately do not set
% this flag: they mirror exactly what Beamer's own \cs{mode*} already
% knows about.}
\cs_new_protected:Npn \osglecture_modes_declare_allowed_command:nn #1#2
  { \prop_gput:Nnn \g_osglecture_modes_dispatch_prop {#1} {#2} }
\cs_generate_variant:Nn \osglecture_modes_declare_allowed_command:nn { xx }

\cs_new_protected:Npn \osglecture_modes_register_command:nnn #1#2#3
  {
    \osglecture_modes_declare_allowed_command:nn {#1} {#2}
    \tl_if_empty:nF {#3}
      { \prop_gput:Nnn \g_osglecture_modes_dispatch_mode_prop {#1} {#3} }
  }
\cs_generate_variant:Nn \osglecture_modes_register_command:nnn { VVV }

\cs_new_protected:Npn \osglecture_modes_keep_generic:Nn #1#2
  {
    #1 {#2}
    \bool_gset_false:N \g_osglecture_modes_after_structure_bool
    \osglecture_modes_scan:
  }

\cs_new_protected:Npn \osglecture_modes_keep_generic_label:Nn #1#2
  {
    \bool_if:NT \g_osglecture_modes_after_structure_bool
      { #1 {#2} }
    \bool_gset_false:N \g_osglecture_modes_after_structure_bool
    \osglecture_modes_scan:
  }

\tl_new:N \l_osglecture_modes_reg_cmd_tl
\tl_new:N \l_osglecture_modes_reg_handler_tl

\NewDocumentCommand \DeclareLectureKeepCommand { O{} m }
  {
    \tl_clear:N \l_osglecture_modes_keep_mode_tl
    \keys_set:nn { osglecture_modes / keep } {#1}
    \bool_gset_true:N \g_osglecture_modes_scan_extended_bool
    \cs_if_exist:NTF #2
      {
        \exp_args:Nc \cs_set_protected:Npn
          { osglecture_modes_keep_generic_ \cs_to_str:N #2 : }
          { \osglecture_modes_keep_generic:Nn #2 }
        \tl_set:Nx \l_osglecture_modes_reg_cmd_tl { \cs_to_str:N #2 }
        \tl_set:Nx \l_osglecture_modes_reg_handler_tl
          { osglecture_modes_keep_generic_ \cs_to_str:N #2 : }
        \osglecture_modes_register_command:VVV
          \l_osglecture_modes_reg_cmd_tl
          \l_osglecture_modes_reg_handler_tl
          \l_osglecture_modes_keep_mode_tl
      }
      { \msg_error:nnN { kernel } { command-not-defined } #2 }
  }

\NewDocumentCommand \DeclareLectureKeepLabelCommand { O{} m }
  {
    \tl_clear:N \l_osglecture_modes_keep_mode_tl
    \keys_set:nn { osglecture_modes / keep } {#1}
    \bool_gset_true:N \g_osglecture_modes_scan_extended_bool
    \cs_if_exist:NTF #2
      {
        \exp_args:Nc \cs_set_protected:Npn
          { osglecture_modes_keep_generic_ \cs_to_str:N #2 : }
          { \osglecture_modes_keep_generic_label:Nn #2 }
        \tl_set:Nx \l_osglecture_modes_reg_cmd_tl { \cs_to_str:N #2 }
        \tl_set:Nx \l_osglecture_modes_reg_handler_tl
          { osglecture_modes_keep_generic_ \cs_to_str:N #2 : }
        \osglecture_modes_register_command:VVV
          \l_osglecture_modes_reg_cmd_tl
          \l_osglecture_modes_reg_handler_tl
          \l_osglecture_modes_keep_mode_tl
      }
      { \msg_error:nnN { kernel } { command-not-defined } #2 }
  }

\cs_new_protected:Npn \osglecture_modes_scan_token:
  {
    \peek_N_type:TF
      { \osglecture_modes_scan_control_sequence:N }
      { \osglecture_modes_discard_one: }
  }

% \ldeen{Modusgebundene Ausnahmen werden hier direkt gegen den normalen
% Modusgraphen geprüft (\cs{osglecture_modes_if_active:n}) statt gegen
% eine eigene, zweite Vererbungsstruktur -- ein Blattmodus mit mehreren
% Filtermodi als Eltern vereinigt deren Ausnahmen also bereits über die
% ohnehin vorhandene transitive Elternhülle.}{Mode-bound exceptions are
% checked here directly against the normal mode graph
% (\cs{osglecture_modes_if_active:n}) rather than against a second,
% separate inheritance structure -- a leaf mode with several filter modes
% as parents therefore already unions their exceptions via the existing
% transitive parent closure.}
% \ldeen{Verwendet eigene, benannte Variablen statt der generischen
% \cs{l\_tmpa\_tl}/\cs{l\_tmpb\_tl}: Der Handlername muss über den Aufruf
% von \cs{osglecture_modes_if_active:n} hinweg erhalten bleiben, das
% \cs{l\_tmpa\_tl} intern selbst für die Aliasauflösung
% verwendet.}{Uses dedicated, named variables instead of the generic
% \cs{l\_tmpa\_tl}/\cs{l\_tmpb\_tl}: the handler name must survive the
% call to \cs{osglecture_modes_if_active:n}, which itself uses
% \cs{l\_tmpa\_tl} internally for alias resolution.}
\tl_new:N \l_osglecture_modes_scan_key_tl
\tl_new:N \l_osglecture_modes_scan_handler_tl
\tl_new:N \l_osglecture_modes_scan_dispatch_mode_tl
\cs_generate_variant:Nn \prop_item:Nn { NV }
\cs_new_protected:Npn \osglecture_modes_scan_control_sequence:N #1
  {
    \tl_set:Nx \l_osglecture_modes_scan_key_tl { \cs_to_str:N #1 }
    \tl_set:Nx \l_osglecture_modes_scan_handler_tl
      {
        \prop_item:NV \g_osglecture_modes_dispatch_prop
          \l_osglecture_modes_scan_key_tl
      }
    \tl_if_empty:NTF \l_osglecture_modes_scan_handler_tl
      { \osglecture_modes_scan: }
      {
        \tl_set:Nx \l_osglecture_modes_scan_dispatch_mode_tl
          {
            \prop_item:NV \g_osglecture_modes_dispatch_mode_prop
              \l_osglecture_modes_scan_key_tl
          }
        \tl_if_empty:NTF \l_osglecture_modes_scan_dispatch_mode_tl
          { \use:c { \l_osglecture_modes_scan_handler_tl } }
          {
            \osglecture_modes_if_active:VTF \l_osglecture_modes_scan_dispatch_mode_tl
              { \use:c { \l_osglecture_modes_scan_handler_tl } }
              { \osglecture_modes_scan: }
          }
      }
  }

\cs_new_protected:Npn \osglecture_modes_keep_part:
  { \osglecture_modes_keep_structure:N \part }
\cs_new_protected:Npn \osglecture_modes_keep_section:
  { \osglecture_modes_keep_structure:N \section }
\cs_new_protected:Npn \osglecture_modes_keep_subsection:
  { \osglecture_modes_keep_structure:N \subsection }
\cs_new_protected:Npn \osglecture_modes_keep_subsubsection:
  { \osglecture_modes_keep_structure:N \subsubsection }
\cs_new_protected:Npn \osglecture_modes_keep_appendix:
  {
    \appendix
    \bool_gset_true:N \g_osglecture_modes_after_structure_bool
    \osglecture_modes_scan:
  }

\NewDocumentCommand \osglecture_modes_keep_lecturemode: { d<> +m }
  {
    \IfNoValueTF {#1}
      { \lecturemode{#2} }
      { \lecturemode<#1>{#2} }
    \bool_gset_false:N \g_osglecture_modes_after_structure_bool
    \osglecture_modes_scan:
  }

\NewDocumentCommand \osglecture_modes_keep_mode: { d<> +m }
  {
    \IfNoValueTF {#1}
      { \mode{#2} }
      { \mode<#1>{#2} }
    \bool_gset_false:N \g_osglecture_modes_after_structure_bool
    \osglecture_modes_scan:
  }

\NewDocumentCommand \osglecture_modes_keep_maketitle: { O{} }
  {
    \bool_gset_false:N \g_osglecture_modes_scanner_active_bool
    \tl_if_blank:nTF {#1} { \maketitle } { \maketitle[#1] }
  }

\osglecture_modes_declare_allowed_command:nn { part }          { osglecture_modes_keep_part: }
\osglecture_modes_declare_allowed_command:nn { section }       { osglecture_modes_keep_section: }
\osglecture_modes_declare_allowed_command:nn { subsection }    { osglecture_modes_keep_subsection: }
\osglecture_modes_declare_allowed_command:nn { subsubsection } { osglecture_modes_keep_subsubsection: }
\osglecture_modes_declare_allowed_command:nn { appendix }      { osglecture_modes_keep_appendix: }
\osglecture_modes_declare_allowed_command:nn { lecturemode }   { osglecture_modes_keep_lecturemode: }
\osglecture_modes_declare_allowed_command:nn { mode }          { osglecture_modes_keep_mode: }
\osglecture_modes_declare_allowed_command:nn { maketitle }     { osglecture_modes_keep_maketitle: }

\NewDocumentCommand \osglecture_modes_keep_structure:N { m s o m }
  {
    \IfBooleanTF {#2}
      {
        \IfNoValueTF {#3}
          { #1*{#4} }
          { #1*[#3]{#4} }
      }
      {
        \IfNoValueTF {#3}
          { #1{#4} }
          { #1[#3]{#4} }
      }
    \bool_gset_true:N \g_osglecture_modes_after_structure_bool
    \osglecture_modes_scan:
  }

\cs_new_protected:Npn \osglecture_modes_discard_one:
  { \osglecture_modes_discard_one:n }
\cs_new_protected:Npn \osglecture_modes_discard_one:n #1
  {
    \bool_gset_false:N \g_osglecture_modes_after_structure_bool
    \osglecture_modes_scan:
  }

\AddToHook{cmd/maketitle/after}
  {
    \bool_gset_true:N \g_osglecture_modes_scanner_active_bool
    \bool_gset_false:N \g_osglecture_modes_after_structure_bool
    \osglecture_modes_start_scan:
  }

% \ldeen{Entscheidet, wie \cs{mode*} bzw. die Schalterform \cs{mode<name>}
% ausgeführt wird. \code{*} delegiert an ein natives \cs{mode*}, wenn
% eines existiert \emph{und} keine eigene Ausnahme über die native
% Positivliste hinausgeht (\cs{g\_osglecture\_modes\_scan\_extended\_bool}) --
% das erhält die native Semantik unverändert im einfachen Fall (siehe den
% \code{native-adapter}-Test). Jeder andere Filtermodus hat naturgemäß
% keine native Entsprechung und läuft daher immer über den eigenen
% Scanner.}{Decides how \cs{mode*} or the switch form \cs{mode<name>} is
% carried out. \code{*} delegates to a native \cs{mode*} when one exists
% \emph{and} no custom exception goes beyond the native allow-list
% (\cs{g\_osglecture\_modes\_scan\_extended\_bool}) -- this keeps native
% semantics unchanged in the simple case (see the \code{native-adapter}
% test). Any other filter mode by construction has no native counterpart
% and therefore always runs through the package's own scanner.}
% \ldeen{Engagiert nur den Scanner selbst (\cs{g\_osglecture\_modes\_scan\_engaged\_bool});
% ob tatsächlich gescannt wird, entscheidet weiterhin
% \cs{g\_osglecture\_modes\_scanner\_active\_bool}, das erst eine
% registrierte Unit-Umgebung oder \cs{maketitle} setzt. Ein früheres
% sofortiges Scharfschalten direkt an dieser Stelle hätte auch
% Präambel-nahen Inhalt gleich nach \cs{begindocument/end} erfasst (etwa
% den Testtreiber-Makro \cs{START} selbst) -- deutlich aggressiver als
% Beamers natives \cs{mode*}, dessen feste Ausnahmeliste genau solche
% frühen Fälle vorsieht.}{Only engages the scanner itself
% (\cs{g\_osglecture\_modes\_scan\_engaged\_bool}); whether scanning
% actually happens is still decided by
% \cs{g\_osglecture\_modes\_scanner\_active\_bool}, which only a
% registered unit environment or \cs{maketitle} sets. Arming scanning
% immediately right here would also have caught preamble-adjacent content
% right after \cs{begindocument/end} (such as the test driver's own
% \cs{START} macro) -- notably more aggressive than Beamer's native
% \cs{mode*}, whose fixed allow-list accounts for exactly such early
% cases.}
\cs_new_protected:Npn \osglecture_modes_switch_own:
  {
    \bool_gset_true:N \g_osglecture_modes_scan_engaged_bool
    \osglecture_modes_start_scan:
  }

\cs_new_protected:Npn \osglecture_modes_switch:n #1
  {
    \osglecture_modes_activate_runtime:n {#1}
    \str_if_eq:nnTF {#1} {*}
      {
        \bool_if:NTF \g_osglecture_modes_native_mode_exists_bool
          {
            \bool_if:NTF \g_osglecture_modes_scan_extended_bool
              { \osglecture_modes_switch_own: }
              { \osglecture_modes_native_mode * }
          }
          { \osglecture_modes_switch_own: }
      }
      { \osglecture_modes_switch_own: }
  }
\cs_generate_variant:Nn \osglecture_modes_switch:n { V }

\AddToHook{begindocument/end}
  {
    \bool_if:NT \g_osglecture_modes_ignorenonframetext_bool
      { \osglecture_modes_switch:n {*} }
  }

% \ldeen{Löst die Kompatibilitätsoption auf, damit Metadaten in der
% Präambel den gewählten Modus sehen. osglecture sollte einen expliziten
% Blattmodus übergeben. Bei \code{deferred}: Eine Klasse, die ihre eigene
% Blattmodus-Auflösung vor \cs{LoadClass} durchführt, überspringt dies
% vollständig.}{Resolve the compatibility option, so metadata in the
% preamble sees the selected mode. osglecture should pass an explicit
% leaf mode. Deferred: a class doing its own leaf resolution before
% \cs{LoadClass} skips this entirely.}
\cs_new_protected:Npn \__osglecture_modes_resolve_leaf:
  {
\tl_if_eq:NnTF \g_osglecture_modes_option_mode_tl {auto}
  {
    \@ifclassloaded{beamer}
      {
        \str_if_eq:eeTF { \use:c { beamer@currentmode } } {handout}
          {
          \SetLectureModes{handout,beamer}
          }
          { \SetLectureModes{slides,beamer} }
      }
      {
        \@ifclassloaded{ltx-talk}
          {
            \cs_if_exist:NTF \l__talk_mode_str
              {
                \str_case:VnF \l__talk_mode_str
                  {
                    { handout }  { \SetLectureModes{handout} }
                    { article }  { \SetLectureModes{article} }
                    { projector }{ \SetLectureModes{slides} }
                  }
                  {
                    \seq_if_in:NVTF \g_osglecture_modes_known_seq \l__talk_mode_str
                      { \exp_args:NV \SetLectureModes \l__talk_mode_str }
                      { \SetLectureModes{slides} }
                  }
              }
              { \SetLectureModes{slides} }
          }
          { \SetLectureModes{article} }
      }
  }
  { \exp_args:NV \SetLectureModes \g_osglecture_modes_option_mode_tl }
  }

% \ldeen{Verdrahtet portable Modusausdrücke mit dem geladenen
% Overlaybackend. Braucht die vorhandene Basisklasse, daher führt ein
% verzögertes Laden dies erst über \cs{LectureModesActivateOverlays} nach
% \cs{LoadClass} aus, nicht sofort. Der Rumpf steckt in einer eigenen
% Datei, damit seine Befehlsparameter-Token nicht eine Verschachtelungsebene
% zu tief innerhalb einer Wrapper-Definition erneut geparst werden.}{Wire
% portable mode expressions into the loaded overlay backend. Needs the
% base class present, so a deferred load runs this from
% \cs{LectureModesActivateOverlays} after \cs{LoadClass} rather than now.
% The body is kept in a separate file so its command-parameter tokens are
% not reparsed one nesting level too deep inside a wrapper definition.}
\cs_new_protected:Npn \__osglecture_modes_patch_overlays:
  { \file_input:n { osglecture-modes-overlays.code.tex } }
\NewDocumentCommand \LectureModesActivateOverlays { }
  { \__osglecture_modes_patch_overlays: }

% \ldeen{Ein gewöhnliches Laden per \cs{usepackage} geschieht nach
% \cs{documentclass}, die Basisklasse ist also bereits vorhanden: sofort
% auflösen und patchen. Ein verzögertes Laden wartet dagegen darauf, dass
% die Klasse \cs{LectureModesActivateOverlays} aufruft.}{A normal
% \cs{usepackage} load happens after \cs{documentclass}, so the base
% class is already present: resolve and patch straight away. A deferred
% load waits for the class to call \cs{LectureModesActivateOverlays}.}
\bool_if:NF \g_osglecture_modes_deferred_bool
  {
    \__osglecture_modes_resolve_leaf:
    \__osglecture_modes_patch_overlays:
  }

\AddToHook{begindocument/before}{\osglecture_modes_finalize_commands:}
%</package>

%<*overlays>
% \ldeen{Wird von \cs{__osglecture_modes_patch_overlays:} in das geladene
% Overlaybackend eingebunden; läuft in der von diesem Aufrufer
% bereitgestellten expl3-Syntax.}{Wired into the loaded overlay backend by
% \cs{__osglecture_modes_patch_overlays:}; runs in expl3 syntax provided
% by its caller.}
% \ldeen{Lässt die nativen Overlay-Decoder einen rein portablen
% Modusausdruck erkennen. Alles andere, insbesondere gemischte
% Modus-/Overlaysyntax, bleibt unangetastet.}{Let the native overlay
% decoders recognize a pure portable mode expression. Everything else,
% especially mixed mode/overlay syntax, remains untouched.}
\bool_if:NT \g_osglecture_modes_patch_overlays_bool
  {
    \cs_if_exist:NT \beamer@masterdecode
      {
        \cs_new_eq:NN \osglecture_modes_native_beamer_masterdecode
          \beamer@masterdecode
        \cs_set_protected:Npn \beamer@masterdecode #1
          {
            \osglecture_modes_classify:n {#1}
            \str_case:Vn \l_osglecture_modes_action_tl
              {
                { show } { \osglecture_modes_native_beamer_masterdecode {all} }
                { hide } { \osglecture_modes_native_beamer_masterdecode {0} }
                { native }
                  {
                    \exp_args:NV \osglecture_modes_native_beamer_masterdecode
                      \l_osglecture_modes_native_tl
                  }
              }
          }
      }
    \cs_if_exist:NT \__talk_decode_parse:n
      {
        \cs_new_eq:NN \osglecture_modes_native_talk_decode_parse:n
          \__talk_decode_parse:n
        \cs_set_protected:Npn \__talk_decode_parse:n #1
          {
            \osglecture_modes_classify:n {#1}
            \str_case:Vn \l_osglecture_modes_action_tl
              {
                { show }
                  { \osglecture_modes_native_talk_decode_parse:n {all} }
                { hide }
                  {
                    \str_if_eq:VnTF \l__talk_mode_str {article}
                      { \osglecture_modes_native_talk_decode_parse:n {projector} }
                      { \osglecture_modes_native_talk_decode_parse:n {article} }
                  }
                { native }
                  {
                    \exp_args:NV \osglecture_modes_native_talk_decode_parse:n
                      \l_osglecture_modes_native_tl
                  }
              }
          }
      }
  }

% \ldeen{Ergänzt ein bestehendes Backend um portable Inlinemodi, ohne
% dessen Umschaltsyntax oder unbekannte/komplexe Overlayspezifikationen
% zu übernehmen.}{Add portable inline modes to an existing backend
% without taking over its switching syntax or unknown/complex overlay
% specifications.}
\cs_if_exist:NTF \mode
  {
    \bool_if:NT \g_osglecture_modes_patch_overlays_bool
      {
        \cs_new_eq:NN \osglecture_modes_native_mode \mode
        \bool_gset_true:N \g_osglecture_modes_native_mode_exists_bool
        % \ldeen{Eine Schalterform \cs{mode<name>} ohne Inhalt wird nur dann
        % als portable Filtermodus-Aktivierung behandelt, wenn \code{name}
        % ein registrierter Modus ist; ein unbekannter (also nativer)
        % Name wird unverändert an das Backend weitergereicht, genau wie
        % zuvor.}{A switch form \cs{mode<name>} without content is only
        % treated as a portable filter-mode activation when \code{name}
        % is a registered mode; an unknown (i.e. native) name is
        % forwarded to the backend unchanged, exactly as before.}
        \cs_new_protected:Npn \osglecture_modes_switch_or_native:n #1
          {
            \osglecture_modes_resolve:nN {#1} \l_tmpa_tl
            \seq_if_in:NVTF \g_osglecture_modes_known_seq \l_tmpa_tl
              { \osglecture_modes_switch:n {#1} }
              { \osglecture_modes_native_mode <#1> }
          }
        \RenewDocumentCommand \mode { s d<> }
          {
            \IfBooleanTF {#1}
              { \osglecture_modes_switch:n {*} }
              {
                \IfNoValueTF {#2}
                  { \osglecture_modes_native_mode }
                  {
                    \peek_meaning:NTF \c_group_begin_token
                      { \osglecture_modes_mode_inline:n {#2} }
                      { \osglecture_modes_switch_or_native:n {#2} }
                  }
              }
          }
        \cs_new_protected:Npn \osglecture_modes_mode_native:nn #1#2
          { \osglecture_modes_native_mode <#1> {#2} }
        \cs_new_protected:Npn \osglecture_modes_mode_inline:n #1#2
          {
            \osglecture_modes_classify:n {#1}
            \str_case:Vn \l_osglecture_modes_action_tl
              {
                { show } {#2}
                { hide } { }
                { native }
                  {
                    \exp_args:NV \osglecture_modes_mode_native:nn
                      \l_osglecture_modes_native_tl {#2}
                  }
              }
          }
      }
  }
  {
    \cs_new_protected:Npn \osglecture_modes_mode_fallback_bare:n #1 {#1}
    \cs_new_protected:Npn \osglecture_modes_mode_fallback_inline:n #1#2
      { \osglecture_modes_if_reveal:nT {#1} {#2} }
    \NewDocumentCommand \mode { s d<> }
      {
        \IfBooleanTF {#1}
          { \osglecture_modes_switch:n {*} }
          {
            \IfNoValueTF {#2}
              {
                \peek_meaning:NTF \c_group_begin_token
                  { \osglecture_modes_mode_fallback_bare:n }
                  { }
              }
              {
                \peek_meaning:NTF \c_group_begin_token
                  { \osglecture_modes_mode_fallback_inline:n {#2} }
                  { \osglecture_modes_switch:n {#2} }
              }
          }
      }
  }

% \ldeen{\cs{only} ist sowohl bei Beamer als auch bei ltx-talk speziell
% implementiert und läuft nicht in jedem Fall durch deren gemeinsame
% Overlay-Decoder, braucht also einen eigenen Adapter statt sich den von
% \cs{mode} zu teilen. Spiegelt die zweigeteilte Form von \cs{mode} oben:
% Existiert \cs{only} bereits (Beamer, ltx-talk), wird es über den
% signaturbewussten Adapter gebunden, sodass portable Modi direkt
% behandelt und native Overlayspezifikationen unverändert weitergereicht
% werden; existiert es nicht (Artikel, scrbook, book, \dots{}), wird ein
% rein portables \cs{only} definiert, damit gemeinsamer Quelltext es
% trotzdem nutzen kann.}{\cs{only} is implemented specially by both
% beamer and ltx-talk and does not pass through their common overlay
% decoders in every case, so it needs its own adapter rather than sharing
% \cs{mode}'s. Mirrors \cs{mode}'s two-branch shape above: where
% \cs{only} already exists (beamer, ltx-talk), bind it through the
% signature-aware adapter so portable modes are handled directly while
% native overlay specifications are delegated unchanged; where it does
% not exist (article, scrbook, book, ...), define a portable-only
% \cs{only} so shared source can still use it.}
\cs_if_exist:NTF \only
  {
    \bool_if:NT \g_osglecture_modes_patch_overlays_bool
      { \MakeOverlayAwareCommand[arguments={+m}]{\only} }
  }
  {
    \NewDocumentCommand \only { d<> +m }
      {
        \IfNoValueTF {#1}
          {#2}
          { \osglecture_modes_if_reveal:nT {#1} {#2} }
      }
  }

% \ldeen{Anders als \cs{only} und \cs{mode} laufen \cs{uncover},
% \cs{visible}, \cs{onslide}, \cs{invisible} und \cs{alt} durch die
% gemeinsamen Overlay-Decoder von Beamer und ltx-talk (oben gepatcht),
% sodass eine portable Modusspezifikation sie dort bereits korrekt
% erreicht, ohne jeden Befehl einzeln anzupassen -- \cs{alt} dient
% zusätzlich als eigenes natives-Spezifikations-Urteil des
% Decoder-Fallbacks weiter unten, ein Einpacken hier würde dieses Urteil
% also auf sich selbst zurückrufen lassen. Nur eine Basisklasse, der der
% Befehl nativ fehlt (Artikel, scrbook, book, \dots{}), braucht eine rein
% portable Definition.}{Unlike \cs{only} and \cs{mode}, \cs{uncover},
% \cs{visible}, \cs{onslide}, \cs{invisible}, and \cs{alt} route through
% beamer's and ltx-talk's common overlay decoders (patched above), so a
% portable mode spec already reaches them correctly there without
% adapting each command individually -- \cs{alt} additionally serves as
% the decoder fallback's own native-spec judge just below, so wrapping it
% here would make that judge call back into itself. Only a base class
% that lacks the command natively (article, scrbook, book, ...) needs a
% portable-only definition.}
\cs_if_exist:NF \uncover
  {
    \NewDocumentCommand \uncover { d<> +m }
      {
        \IfNoValueTF {#1}
          {#2}
          { \osglecture_modes_if_reveal:nT {#1} {#2} }
      }
  }
\cs_if_exist:NF \visible
  {
    \NewDocumentCommand \visible { d<> +m }
      {
        \IfNoValueTF {#1}
          {#2}
          { \osglecture_modes_if_reveal:nT {#1} {#2} }
      }
  }
\cs_if_exist:NF \onslide
  {
    \NewDocumentCommand \onslide { d<> +m }
      {
        \IfNoValueTF {#1}
          {#2}
          { \osglecture_modes_if_reveal:nT {#1} {#2} }
      }
  }

% \ldeen{\cs{invisible} kehrt die Polarität um: Inhalt wird überall
% angezeigt außer im angegebenen Modus. Eine native oder unerkannte
% Spezifikation degradiert wie bei den positiven Befehlen zu
% \emph{anzeigen}.}{\cs{invisible} inverts the polarity: content shows
% everywhere except in the given mode. A native or unrecognised
% specification degrades to "show" like the positive commands.}
\cs_if_exist:NF \invisible
  {
    \NewDocumentCommand \invisible { d<> +m }
      {
        \IfNoValueTF {#1}
          {#2}
          {
            \osglecture_modes_classify:n {#1}
            \str_if_eq:VnF \l_osglecture_modes_action_tl { show } {#2}
          }
      }
  }

% \ldeen{\cs{alt} nimmt zwei Inhaltsargumente (aktiver Zweig, inaktiver
% Zweig) statt des einen von \cs{only}; der Fallback spiegelt diese Form
% und behandelt eine native oder unerkannte Spezifikation als aktiven
% Zweig.}{\cs{alt} takes two content arguments (active branch, inactive
% branch) instead of \cs{only}'s single one; the fallback mirrors that
% shape and treats a native or unrecognised specification as the active
% branch.}
\cs_if_exist:NF \alt
  {
    \NewDocumentCommand \alt { d<> +m +m }
      {
        \IfNoValueTF {#1}
          {#2}
          {
            \osglecture_modes_classify:n {#1}
            \str_if_eq:VnTF \l_osglecture_modes_action_tl { hide } {#3} {#2}
          }
      }
  }

% \ldeen{\cs{pause} markiert einen Live-Aufdeckpunkt ohne zu wählenden
% Inhalt, und die Vorher-/Während-/Nachher-Aufteilung von \cs{temporal}
% setzt einen geordneten Overlaybereich voraus, den benannte portable
% Modi nicht besitzen -- keines von beidem lässt sich auf
% Inhaltsauswahl per Modus abbilden wie bei den Befehlen oben.
% Basisklassen ohne eigenes \cs{pause}/\cs{temporal} erhalten dennoch eine
% Definition, damit gemeinsamer Quelltext kompiliert: \cs{pause} tut
% nichts, und \cs{temporal} zeigt stets sein \emph{während}-Argument, den
% Zustand, auf den eine statische, overlaylose Darstellung sich
% einpendelt.}{\cs{pause} marks a live-reveal point with no content to
% select, and \cs{temporal}'s before/during/after split assumes an
% ordered overlay range, which named portable modes do not have --
% neither maps onto content selection by mode the way the commands above
% do. Base classes without their own \cs{pause}/\cs{temporal} still get a
% definition so shared source compiles: \cs{pause} does nothing, and
% \cs{temporal} always shows its "during" argument, the state a
% static, non-overlay rendering settles on.}
\cs_if_exist:NF \pause
  { \NewDocumentCommand \pause { O{} } { } }
\cs_if_exist:NF \temporal
  { \NewDocumentCommand \temporal { d<> +m +m +m } {#3} }
%</overlays>

%<*ltxtalk>
\ProvidesExplPackage{osglecture-modes-ltxtalk}
{2026-08-02}
{v0.4.0}
{ltx-talk adapter: lecture-specific baseline exceptions}

\RequirePackage{osglecture-modes}

\ExplSyntaxOn

% \ldeen{Die generische "`außerhalb von Unit-Umgebungen verwerfen"'-Maschinerie
% (Scanner, Modus \code{*}, \cs{DeclareLectureKeepCommand} und Geschwister)
% gehört jetzt zum Kernpaket und funktioniert unabhängig von ltx-talk. Dieser
% Adapter ergänzt nur noch das, was tatsächlich ltx-talk-spezifisch ist:
% \cs{lecture} als weitere Baseline-Ausnahme (mit eigener, dreiteiliger
% Signatur) sowie den Hook, der den Scanner nach \cs{lecture} scharf
% schaltet -- \cs{frame}/\cs{frame*} werden bereits vom Kernpaket erkannt und
% registriert, sobald sie existieren.}{The generic "drop everything outside
% unit environments" machinery (scanner, mode \code{*},
% \cs{DeclareLectureKeepCommand} and siblings) now lives in the core package
% and works independently of ltx-talk. This adapter only adds what is
% actually ltx-talk-specific: \cs{lecture} as a further baseline exception
% (with its own three-part signature) and the hook that arms the scanner
% after \cs{lecture} -- \cs{frame}/\cs{frame*} are already detected and
% registered by the core package once they exist.}

\keys_define:nn { osglecture_modes_ltxtalk }
  {
    mode .code:n =
      { \SetLectureModes{#1} },

    ignorenonframetext .bool_gset:N =
      \g_osglecture_modes_ignorenonframetext_bool,

    ignorenonframetext .default:n = true,

    raw-skip-envs .clist_gset:N =
      \g_osglecture_modes_raw_skip_envs_clist,
  }

\ProcessKeyOptions [ osglecture_modes_ltxtalk ]

\sys_if_engine_luatex:F
  {
    \bool_if:NT \g_osglecture_modes_ignorenonframetext_bool
      {
        \PackageWarning
          {osglecture_modes_ltxtalk}
          {ignorenonframetext with raw-skip support needs LuaLaTeX}
      }
  }

\NewDocumentCommand \osglecture_modes_keep_lecture: { O{} m m }
  {
    \lecture[#1]{#2}{#3}
    \bool_gset_false:N \g_osglecture_modes_after_structure_bool
    \osglecture_modes_scan:
  }
\osglecture_modes_declare_allowed_command:nn { lecture } { osglecture_modes_keep_lecture: }

% \ldeen{\cs{label}, \cs{input} und \cs{include} sind schlichte Instanzen
% der allgemeinen Ausnahmeschnittstelle des Kernpakets; sie über diese statt
% von Hand zu deklarieren, erhält eine einzige Implementierung für
% \emph{diesen Befehl erhalten, dann den Scan fortsetzen}.}{\cs{label},
% \cs{input}, and \cs{include} are plain instances of the core package's
% general exception interface; declaring them through it rather than by
% hand keeps a single implementation for "keep this command, then resume
% scanning".}
\DeclareLectureKeepLabelCommand \label
\DeclareLectureKeepCommand \input
\DeclareLectureKeepCommand \include

% \ldeen{Rüstet den Scanner auf, nachdem osglecture seine
% Gliederungseinheit-Deklaration vollständig verarbeitet hat. Das deckt
% auch Nichtrahmenmaterial vor dem Titel ab.}{Arm the scanner after
% osglecture has completely processed its logical-unit declaration. This
% covers non-frame material before the title as well.}
\AddToHook{cmd/lecture/after}
  {
    \bool_gset_true:N \g_osglecture_modes_scanner_active_bool
    \osglecture_modes_start_scan:
  }

%</ltxtalk>

%<*lua>
osglecture_modes_ltxtalk = osglecture_modes_ltxtalk or {}

osglecture_modes_ltxtalk.raw_skip_env = nil

local previous_callback = callback.find("process_input_buffer")

local function escape_pattern(s)
  return (s:gsub("([^%w])", "%%%1"))
end

function osglecture_modes_ltxtalk.start_raw_skip(env)
  osglecture_modes_ltxtalk.raw_skip_env = env
end

local function raw_skip_line(line)
  local env = osglecture_modes_ltxtalk.raw_skip_env
  if not env then
    return nil
  end

  local pattern =
    "^%s*\\end%s*{%s*" .. escape_pattern(env) .. "%s*}"

  if line:match(pattern) then
    osglecture_modes_ltxtalk.raw_skip_env = nil
  end

  return ""
end

local function process_input_buffer(line)
  local replacement = raw_skip_line(line)

  if replacement ~= nil then
    return replacement
  end

  if previous_callback then
    return previous_callback(line)
  end

  return nil
end

luatexbase.add_to_callback(
  "process_input_buffer",
  process_input_buffer,
  "osglecture-modes-ltxtalk raw skip"
)
%</lua>
